\documentclass[11pt]{article}
% Diplomatic transcription macros (12 Nov 2026). Typed text is plain; the handwritten layer is blue,
% deletions red, typewriter strike-outs grey, doubtful readings in brackets with a question mark.
\usepackage[utf8]{inputenc}
\usepackage[T1]{fontenc}
\usepackage{lmodern}
\usepackage[a4paper,margin=2.2cm]{geometry}
\usepackage{amsmath,amssymb}
\usepackage{xcolor}
\usepackage{tikz-cd}
\usepackage{environ}
\setlength{\parindent}{0pt}
\setlength{\parskip}{4pt}
\title{Notes of the Montreal seminar, 1970 (cote 129, typed part)\\[4pt]\large Diplomatic transcription from the Montpellier scans}
\author{AI-transcribed by Claude Opus 5.5}
\date{12 November 2026}
\newcommand{\scanpage}[1]{\par\bigskip\noindent{\small\sffamily[scan #1]}\par}
\newcommand{\mspage}[1]{{\small\sffamily[p.\,#1]}}
\newcommand{\hand}[1]{\textcolor{blue}{#1}}
\let\over\relax
\newcommand{\over}[2]{\textcolor{gray}{#1}\textcolor{blue}{\ensuremath{^{\!\nearrow}}#2}}
\newcommand{\inserted}[2][]{\textcolor{blue}{$\langle$#2$\rangle$}}
\newcommand{\margin}[1]{\textcolor{blue}{\{\textsf{marg.:} #1\}}}
\newcommand{\deleted}[1]{\textcolor{red}{[\textsf{del.:} #1]}}
\newcommand{\xout}[1]{\textcolor{gray}{[\textsf{x:} #1]}}
\newcommand{\uline}[1]{\underline{#1}}
\newcommand{\uuline}[1]{\underline{\underline{#1}}}
\newcommand{\unsure}[1]{[#1?]}
\newcommand{\illegible}[1]{[\textsf{illeg.\ #1}]}
\newcommand{\diagram}[1]{\par{\centering #1\par}}
\newenvironment{diagramblock}{\par\begin{center}}{\end{center}\par}

\begin{document}
\maketitle

% Notes of the Montreal seminar, 1970 (cote 129, typed part): archive cote 129, site page 1 (PDF page 2).
% typescript. Draft compared: none.
% AI-transcribed by Claude Opus 5.5 (claude-opus-5-5, adaptive thinking, effort high), Message Batches API, batch mtl129_full (msgbatch_01AGnDKHwy6ATCM7ehUa6263), 2026-11-13; not yet checked by eye.
\scanpage{129/1}
\margin{DEL\unsure{A}LE}

CHAPITRE I

PRELIMINAIRES SUR WITT

\mspage{\hand{1}}

\clearpage
% Notes of the Montreal seminar, 1970 (cote 129, typed part): archive cote 129, site page 2 (PDF page 3).
% typescript. Draft compared: none.
% AI-transcribed by Claude Opus 5.5 (claude-opus-5-5, adaptive thinking, effort high), Message Batches API, batch mtl129_full (msgbatch_01AGnDKHwy6ATCM7ehUa6263), 2026-11-13; not yet checked by eye.
\scanpage{129/2}
1.  \uline{Rappels sur les vecteurs de Witt}.

Dans tout ce qui suit, $p$  désigne un nombre premier fixé
une fois pour toutes.

Le \uline{schéma des vecteurs de Witt} $\mathbb{W}$  et le \uline{schéma des vecteurs}
\uline{de Witt tronqués à l'ordre  $n$} $\mathbb{W}_n$  sont des schémas en anneaux, affines
sur $\mathbb{Z}$ qui sont définis comme suit:

Soit  $n \geq 1$  un entier et soit  $\mathbb{E}^n = \mathrm{Spec}\, \mathbb{Z}[T_1,T_2,\dots,T_n]$
d'espace affine de dimension  $n$  sur  $\mathbb{Z}$.  Le schéma  $\mathbb{E}^n$  représente
le foncteur  $A \mapsto A^n$  de la catégorie des schémas affines dans la
catégorie des ensembles; il est donc muni d'une structure naturelle de
schéma en anneaux.  On note  $\underline{O}^n$  le schéma  $\mathbb{E}^n$  muni de cette structure.
On démontre alors (c.f. J.-P. Serre, ``Corps locaux'' II §6) qu'il existe
une unique structure de schéma en anneaux sur  $\mathbb{E}^n$  telle que le morphisme
de schémas
\[
\phi = (\phi_1,\phi_2,\dots,\phi_n) \colon \mathbb{E}^n \to \underline{O}^n
\]
défini par
\[
\phi_i(x_1,x_2,\dots,x_n) = x_1^{p^{i-1}} + p\, x_2^{p^{i-2}} + p^2\, x_3^{p^{i-3}} + \dots + p^{i-1}\, x_i
\]
soit un homomorphisme d'anneaux.

Par définition, $\mathbb{W}_n$  est le schéma  $\mathbb{E}_n$  muni de cette structure
de schéma en anneaux.

\mspage{\hand{2}}

\clearpage
% Notes of the Montreal seminar, 1970 (cote 129, typed part): archive cote 129, site page 3 (PDF page 4).
% typescript. Draft compared: none.
% AI-transcribed by Claude Opus 5.5 (claude-opus-5-5, adaptive thinking, effort high), Message Batches API, batch mtl129_full (msgbatch_01AGnDKHwy6ATCM7ehUa6263), 2026-11-13; not yet checked by eye.
\scanpage{129/3}
Les morphismes de restriction
\[ R_n : \mathbb{W}_{n+1} \to \mathbb{W}_n \qquad (x_1,\ldots,x_n,x_{n+1}) \mapsto (x_1,\ldots,x_n) \]
sont alors des homomorphismes d'anneaux. En effet, il résulte de
l'unicité de la structure d'anneau sur $\mathbb{W}_n$ qu'un morphisme $f: X \to \mathbb{W}_n$
est un morphisme d'anneaux (resp. de groupes) si et seulement si
$\phi_i \circ f: X \to \underline{O}$ est un homomorphisme d'anneaux (resp. de groupes pour
tout $i$, $1 \leq i \leq n$. Ici, on a $\phi_i\ R_n = \phi_i$, qui est par définition
un homomorphisme d'anneaux.

On définit le schéma en anneau des vecteurs de Witt $\mathbb{W}$ par
\[ \mathbb{W} = \lim_{\substack{\leftarrow \\ n}} (\mathbb{W}_n,R_n) \]
Le schéma sous-jacent à $\mathbb{W}$ est donc en particulier isomorphe à $\mathbb{E}^{\mathbb{N}}$.

2. \uline{Quelques morphismes remarquables}.

2.1. On définit
\[ V = V_n : \mathbb{W}_n \to \mathbb{W}_n \qquad (x_1,x_2,\ldots,x_{n-1},x_n) \mapsto (0,x_1,\ldots,x_{n-1}) \]
\[ T = T_n : \mathbb{W}_n \to \mathbb{W}_{n+1} \qquad (x_1,x_2,\ldots,x_n) \mapsto (0,x_1,\ldots,x_n). \]
Le morphisme $T$ est \uline{additif}; on a en effet $\phi_i\ T = p\phi_{i-1}$
pour $1 \leq i \leq n+1$.

On en déduit que $V_n = R_n\ T_n = T_{n-1}\ R_{n-1}$ est aussi un morphisme
\uline{additif}.

\mspage{\hand{3}}

\clearpage
% Notes of the Montreal seminar, 1970 (cote 129, typed part): archive cote 129, site page 4 (PDF page 5).
% typescript. Draft compared: none.
% AI-transcribed by Claude Opus 5.5 (claude-opus-5-5, adaptive thinking, effort high), Message Batches API, batch mtl129_full (msgbatch_01AGnDKHwy6ATCM7ehUa6263), 2026-11-13; not yet checked by eye.
\scanpage{129/4}
Par passage à la limite, les morphismes $V_n$ définissent
un morphisme additif
\[ V\colon \mathbb{W} \to \mathbb{W} \]
appelé \uline{Verschiebung} ou \uline{morphisme de décalage}.

2.2. Le \uline{morphisme de frobénius} $F\colon \mathbb{W} \to \mathbb{W}$ est défini par passage
à la limite à partir des morphismes
\[ F_n\colon \mathbb{W}_n \to \mathbb{W}_n \qquad (x_1,x_2,\dots,x_n) \mapsto (x_1^p,x_2^p,\dots,x_n^p). \]

Ce morphisme n'est en général par additif, mais par réduction
en caractéristique $p$, il définit un \uline{homomorphisme d'Anneau}
\[ F_{\mathbb{F}_p} : \mathbb{W}_{\mathbb{F}_p} \to \mathbb{W}_{\mathbb{F}_p} \qquad (\text{resp. } F_{n\,\mathbb{F}_p}) \]
où $\mathbb{F}_p$ désigne le corps premier $\mathbb{Z}/p\mathbb{Z}$ et $\mathbb{W}_{\mathbb{F}_p}$ désigne le schéma
induit sur $\mathbb{F}_p$ par changement de base. En effet les deux lois de
composition sur $\mathbb{W}$ sont données par
\[ (x_1,x_2,\dots,x_n,\dots) + (y_1,\dots,y_n,\dots) = (S_1,S_2,\dots,S_n,\dots) \]
\[ (x_1,x_2,\dots,x_n,\dots)\cdot(y_1,\dots,y_n,\dots) = (P_1,\dots,P_n,\dots) \]
où les $S_i$ et $P_i$ sont des polynômes à coefficients entiers en
$(x_1,x_2,\dots,x_i,\ y_1,y_2,\dots,y_i)$. Or pour toute $\mathbb{F}_p$-algèbre, on a
\[ P(x^p,y^p) = [P(x,y)]^p. \]

\mspage{\hand{4}}

\clearpage
% Notes of the Montreal seminar, 1970 (cote 129, typed part): archive cote 129, site page 5 (PDF page 6).
% typescript. Draft compared: none.
% AI-transcribed by Claude Opus 5.5 (claude-opus-5-5, adaptive thinking, effort high), Message Batches API, batch mtl129_full (msgbatch_01AGnDKHwy6ATCM7ehUa6263), 2026-11-13; not yet checked by eye.
\scanpage{129/5}
Proposition (2.3). En caractéristique p, on a les relations
\[ F.V = V.F = p\ \mathrm{Id}_{\mathbb{W}_{\mathbb{F}_p}} \]
(c.f. J.-P. Serre "corps locaux" II §6 Cor. du th.7). On peut
donner la démonstration directe suivante: Posons
\[ h = (h_1,h_2,\ldots,h_n,\ldots) = (p.\mathrm{Id}_W - VF)(x_1,\ldots,x_n,\ldots). \]
On sait que $h_i \in \mathbb{Z}[x_1,x_2,\ldots,x_i]$, et il suffit de démontrer que $h_i$
est divisible par $p$ pour tout $i$ pour prouver la proposition. Or
on trouve que
\[ \phi_i(h) = p\phi_i(x) - \phi_i(VF(x)) = p^i\, x_i \]
d'où l'on déduit que $h_1 = px_1$,
\[ \phi_2(h) = h_1^p + ph_2 = p^2\, x_2 \quad \text{d'où} \quad h_2 = p(x_2-p^{p-2}\, x_1), \]
et par récurrence sur $i$, on en déduit que pour tout $i$ on a $h_i = pk_i$.

Notons en particulier que dans $\mathbb{W}_{\mathbb{F}_p}$ on a $p = (0,1,0,\ldots)$
(car $(1,0,0,\ldots)$ est l'unité de $\mathbb{W}$).

3. \uline{L'ind-schéma} $\underset{\rightarrow}{\mathbb{W}}$ \uline{et l'ind-schéma} $\underset{\rightarrow}{\mathbb{W}}'$

On définit
\[ \underset{\rightarrow}{\mathbb{W}} = \underset{n}{\underset{\rightarrow}{\lim}}\, (\mathbb{W}_n,T_n). \]
C'est un ind-schéma en groupes dont le ind-schéma sous-jacent est

\mspage{\unsure{5}}

\clearpage
% Notes of the Montreal seminar, 1970 (cote 129, typed part): archive cote 129, site page 6 (PDF page 7).
% typescript. Draft compared: none.
% AI-transcribed by Claude Opus 5.5 (claude-opus-5-5, adaptive thinking, effort high), Message Batches API, batch mtl129_full (msgbatch_01AGnDKHwy6ATCM7ehUa6263), 2026-11-13; not yet checked by eye.
\scanpage{129/6}
isomorphe à  $\mathbb{E}^{(\mathbb{N})}$, c'est-à-dire que pour tout anneau  A,
\[ \underset{\rightarrow}{\mathbb{W}}(A) = A^{(\mathbb{N})} = \{(x_1,x_2,\ldots,x_n,\ldots),\ x_i \in A,\ x_i = 0 \text{ sauf un nombre fini}\}. \]
Mais ce foncteur a une structure additive qui n'est pas la structure
habituelle de $\mathbb{E}^{(\mathbb{N})}$.

On peut définir un autre système inductif naturel à l'aide des
morphismes
\[ \mathbb{W}_n \xrightarrow{T_n F_n} \mathbb{W}_{n+1} . \]
On a alors des diagrammes commutatifs
\begin{diagramblock}
\begin{tikzcd}
\mathbb{W}_n \arrow[r, "T_n"] \arrow[d, "F_n^{n-1}"] & \mathbb{W}_{n+1} \arrow[d, "F_{n+1}^{n}"] \\
\mathbb{W}_n \arrow[r, "T_nF_n = p"] & \mathbb{W}_{n+1}
\end{tikzcd}
\end{diagramblock} ,
car  $T_nF_n = F_{n+1}\,T_n$.  Par définition, on pose
\[ \underset{\rightarrow}{\mathbb{W}}' = \underset{\rightarrow}{\lim}\ (\mathbb{W}_n,\ T_nF_n). \]
En caractéristique  0, $\underset{\rightarrow}{\mathbb{W}}'$  n'a aucune structure naturelle de Groupe,
mais en caractéristique p $\underset{\rightarrow}{\mathbb{W}}'$ est un ind-schéma en groupes, et on a un
morphisme additif
\[ u\colon \underset{\rightarrow}{\mathbb{W}}_{\mathbb{F}_p} \to \underset{\rightarrow}{\mathbb{W}}'_{\mathbb{F}_p} \]
obtenu par passage à la limite des  $F^{n-1}$.

\mspage{\hand{6}}

\clearpage
% Notes of the Montreal seminar, 1970 (cote 129, typed part): archive cote 129, site page 7 (PDF page 8).
% typescript. Draft compared: none.
% AI-transcribed by Claude Opus 5.5 (claude-opus-5-5, adaptive thinking, effort high), Message Batches API, batch mtl129_full (msgbatch_01AGnDKHwy6ATCM7ehUa6263), 2026-11-13; not yet checked by eye.
\scanpage{129/7}
Notons qu'en caractéristique $p$, les morphismes de transition $TF$ ne
sont autres que \uline{la multiplication par $p$} puisque l'on a $TFR = FTR = FV = p.Id$.

\uline{Remarque}. Si $A$ est un anneau \uline{parfait} de caractéristique $p$, $u$ induit
un isomorphisme
\[ u(A) : \mathbb{W}_{\to\mathbb{F}_p}(A) \to \mathbb{W}'_{\to\mathbb{F}_p}(A) \]
puisque c'est vrai pour chaque $F^n(A)$.

4. \uline{L'anneau} $W$.

4.1. On rappelle que (voir J.P. Serre ``corps locaux'' II §6 th.7)
que si $k$ est un corps parfait de caractéristique $p$, l'anneau
\[ W = \mathbb{W}(k) = \lim_{\leftarrow} W_n \qquad (W_n = \mathbb{W}_n(k)) \]
est un anneau de valuation discrète, complet, de corps résiduel $k$,
d'uniformisante $p = (0,1,0,0,\ldots)$, dont l'homomorphisme d'augmentation
est donné par la projection canonique
\[ W \to W_1 = k . \]
On peut d'ailleurs montrer (loc. cit. II §5 th.3) que ces propriétés
caractérisent l'anneau $W$ à un isomorphisme unique près (théorème de
Cohen).

De plus, le corps des fractions de $W$,
\[ K = W[1/p] \]
est de caractéristique nulle.

\mspage{\hand{7}}

\clearpage
% Notes of the Montreal seminar, 1970 (cote 129, typed part): archive cote 129, site page 8 (PDF page 9).
% typescript. Draft compared: none.
% AI-transcribed by Claude Opus 5.5 (claude-opus-5-5, adaptive thinking, effort high), Message Batches API, batch mtl129_full (msgbatch_01AGnDKHwy6ATCM7ehUa6263), 2026-11-13; not yet checked by eye.
\scanpage{129/8}
Comme $k$ est parfait, l'homomorphisme $u$ définit un isomor-%
phisme de groupes
\[ u(k) : \underset{\to}{\mathbb{W}}(k) \to \underset{\to}{\mathbb{W}}'(k) . \]

Remarquons alors que la projection naturelle $\pi_n : W \to W_n$
définit un isomorphisme
\[ \varepsilon_n : W/p^nW \xrightarrow{\ \sim\ } W_n \]
et l'on vérifie immédiatement que l'on a un diagramme commutatif
\begin{diagramblock}
\begin{tikzcd}
W \arrow[rrr, "p\,=\,VF"] \arrow[dr] \arrow[dd, "\pi_n"'] & & & W \arrow[dl] \arrow[dd, "\pi_{n+1}"] \\
& W/p^nW \arrow[r, "p"] \arrow[dl, "\varepsilon_n", "\sim"'] & W/p^{n+1}W \arrow[dr, "\varepsilon_{n+1}"', "\sim"] & \\
W_n \arrow[rrr, "TF\,=\,FT"] & & & W_{n+1}
\end{tikzcd}
\end{diagramblock}

Utilisant alors l'isomorphisme naturel (de $W$-modules) valable pour tout
anneau de valuation discrète
\[ \underset{\to}{\lim}(W/p^nW,p) \xrightarrow{\ \sim\ } K/W \]
on déduit de ce qui précède que $\underset{\to}{\mathbb{W}}'(k)$ et donc $\underset{\to}{\mathbb{W}}(k)$ est isomorphe
au module dualisant $K/W$. Comme $K/W$ est un $W$-module, on déduit une
structure de $W$-module sur $\underset{\to}{\mathbb{W}}(k)$ par transport de structure.

\uline{Proposition} (4.2). Cette structure se prolonge de manière unique en
une structure de ind-schéma en groupes à anneau d'opérateurs $W$ sur
\[ \underset{\to}{\mathbb{W}}_k = \underset{\to}{\lim}((\mathbb{W}_n)_k , T) \]

\mspage{\hand{8}}

\clearpage
% Notes of the Montreal seminar, 1970 (cote 129, typed part): archive cote 129, site page 9 (PDF page 10).
% typescript. Draft compared: none.
% AI-transcribed by Claude Opus 5.5 (claude-opus-5-5, adaptive thinking, effort high), Message Batches API, batch mtl129_full (msgbatch_01AGnDKHwy6ATCM7ehUa6263), 2026-11-13; not yet checked by eye.
\scanpage{129/9}
en faisant opérer $\lambda \in W$ sur $(\mathbb{W}_n)_k$ par multiplication par $F^{1-n}(\lambda)$.

Remarquons que l'élément $F^{1-n}(\lambda) \in W$ est bien défini car,
$k$ étant parfait, $F\colon W \to W$ est un isomorphisme. Cette multiplication
a un sens car, pour toute $k$-algèbre $A$, $\mathbb{W}_n(A)$ est une $W$-algèbre via le
morphisme
\[ W = \mathbb{W}(k) \xrightarrow{\ \pi_n\ } \mathbb{W}_n(k) \to \mathbb{W}_n(A) . \]

Pour que cette action se transporte à $\underrightarrow{\mathbb{W}}_k$, il faut montrer que
l'action de $\lambda \in W$ commute aux morphismes $T$. On vérifie immédiatement
qu'il suffit de montrer que pour toute $k$-algèbre $A$, on a:
\[ \forall\, x \in \mathbb{W}_n(A) ,\ \forall\, \lambda' \in \mathbb{W}_{n+1}(A) \quad \lambda'.T(x) = T(FR(\lambda').x) \]
ou, ce qui revient au même puisque $R$ est une surjection:
\[ \forall\, \lambda' ,\ y \in \mathbb{W}_{n+1}(A) ,\ \ \lambda'.V(y) = V(F(\lambda'),y) . \]
Or si $A$ est un anneau parfait, $F$ est un isomorphisme, et cette relation
est vérifiée car
\[ F(\lambda'.V(y)) = F(\lambda')\ FV(y) = F(\lambda').py = p[F(\lambda').y] = FV[F(\lambda').y]. \]
Dans le cas général, il suffit de remarquer que $\lambda'.V(y)$ et $V(F(\lambda').y)$
s'expriment à l'aide de polynômes à coefficients dans $\mathbb{Z}/p\mathbb{Z}$. D'après
ce qui précède, ces polynômes prennent la même valeur si les variables
sont prises dans un anneau parfait, par exemple la clôture algébrique
de $k$. On en déduit qu'ils sont égaux.

\mspage{\hand{9}}

\clearpage
% Notes of the Montreal seminar, 1970 (cote 129, typed part): archive cote 129, site page 10 (PDF page 11).
% typescript. Draft compared: none.
% AI-transcribed by Claude Opus 5.5 (claude-opus-5-5, adaptive thinking, effort high), Message Batches API, batch mtl129_full (msgbatch_01AGnDKHwy6ATCM7ehUa6263), 2026-11-13; not yet checked by eye.
\scanpage{129/10}
Enfin, il faut montrer que cette structure sur $\underset{\to}{\mathbb{W}}_k$ induit
sur $\underset{\to}{\mathbb{W}}(k)$ la structure précédemment définie. Or il est immédiat
(et cela prouve l'unicité de la structure sur $\underset{\to}{\mathbb{W}}_k$) que la multiplication
par $F^{1-n}(\lambda)$ sur $W_n(k) = W_n$ induit sur $W/p^nW$ la multiplication par
$\lambda$ via les isomorphismes
\[ W_n \xrightarrow[\sim]{F^{n-1}} W_n \xleftarrow[\sim]{\varepsilon_n} W/p^nW \ . \]

\uline{Remarque}. Il n'y a pas en général d'action de $\mathbb{W}_k$ sur $\underset{\to}{\mathbb{W}}_k$ mais seu-%
lement de $W = \mathbb{W}(k)$ sur $\underset{\to}{\mathbb{W}}_k$.

Cette construction est fonctorielle en $k$ et est valable pour
tout anneau parfait $A$ (mais $\mathbb{W}(A)$ n'est plus en général un anneau de
valuation discrète).

5. \uline{Les morphismes de Frobenius et de Verschiebung d'un schéma en groupes}.

Dans tout ce qui suit, on suppose que tous les schémas sont des
schémas sur le corps premier $\mathbb{F}_p$.

5.1. Si $X$ est un schéma, on appelle \uline{morphisme de Frobenius absolu}
de $X$ et on note
\[ \mathbf{f}_X : X \to X \]
l'endomorphisme de $X$ qui est l'identité sur l'espace sous-jacent et
qui associe à une section $s$ de $\underline{\mathcal{O}}_X$ la section $s^p$. Si $Y$ est un
autre schéma, et $g: X \rightharpoonup Y$ un morphisme, le carré suivant est commutatif:

\mspage{\hand{10}}

\clearpage
% Notes of the Montreal seminar, 1970 (cote 129, typed part): archive cote 129, site page 11 (PDF page 12).
% typescript. Draft compared: none.
% AI-transcribed by Claude Opus 5.5 (claude-opus-5-5, adaptive thinking, effort high), Message Batches API, batch mtl129_full (msgbatch_01AGnDKHwy6ATCM7ehUa6263), 2026-11-13; not yet checked by eye.
\scanpage{129/11}
\begin{diagramblock}\begin{tikzcd}
X \arrow[r, "\mathrm{ff}_X"] \arrow[d, "g"'] & X \arrow[d, "g"'] \\
Y \arrow[r, "\mathrm{ff}_Y"] & Y
\end{tikzcd}\end{diagramblock}

On obtient donc un endomorphisme $\mathrm{ff}$ du foncteur identique de $\mathrm{Sch}/\mathbb{F}_p$.

5.2. Soit $S$ un schéma fixé et $X$ un S-schéma. On note $X^{(p/S)}$
ou simplement $X^{(p)}$ le S-schéma image réciproque de $X$ par le changement
de base $\mathrm{ff}_S : S \to S$.

Le diagramme ci-dessus étant commutatif, il existe une flèche
unique $\mathrm{ff}_{X/S}$ telle que le diagramme suivant soit commutatif:

\begin{diagramblock}\begin{tikzcd}
X \arrow[dr, "\mathrm{ff}_{X/S}"'] \arrow[drr, "\mathrm{ff}_X"] \arrow[ddr] & & \\
& X^{(p)} \arrow[r] \arrow[d] & X \arrow[d] \\
& S \arrow[r, "\mathrm{ff}_S"] & S
\end{tikzcd}\end{diagramblock}

On appellera le morphisme $\mathrm{ff}_{X/S}$ ainsi défini le \uline{morphisme de Frobenius}
\uline{relatif de} $X$ \uline{au-dessus de} $S$.

L'application $X \mapsto \mathrm{ff}_{X/S}$ est un morphisme fonctoriel du foncteur
identique de $\mathrm{Sch}/S$ dans le foncteur $X \mapsto X^{(p/S)}$. Ce dernier foncteur
commute aux limites projectives finies, donc si $X$ est un S-groupe,
$X^{(p/S)}$ l'est aussi et, $\mathrm{ff}_{X/S}$ étant fonctoriel, est un morphisme de groupes.

La construction de $X^{(p/S)}$ commute aux changements de base,
c'est-à-dire que pour tout morphisme $T \to S$, on a

\mspage{\hand{11}}

\clearpage
% Notes of the Montreal seminar, 1970 (cote 129, typed part): archive cote 129, site page 12 (PDF page 13).
% typescript. Draft compared: none.
% AI-transcribed by Claude Opus 5.5 (claude-opus-5-5, adaptive thinking, effort high), Message Batches API, batch mtl129_full (msgbatch_01AGnDKHwy6ATCM7ehUa6263), 2026-11-13; not yet checked by eye.
\scanpage{129/12}
\[ (X \underset{S}{\text{x}} T)^{(p/T)} \cong X^{(p/S)} \underset{S}{\text{x}} T \ . \]
Enfin, si $S$ est le spectre du corps premier $\mathbb{F}_p$, $X^{(p/S)} = X$ et
$\mathbb{f}_{X/S} = \mathbb{f}_X$.

5.3. Soit $G$ un S-schéma en groupes commutatifs et plat variable.
On sait définir un homomorphisme canonique de groupes, fonctoriel en $G$
\[ \mathbb{V}_G \colon G^{(p/S)} \to G \]
appelé \uline{homomorphisme de Verschiebung de} $G$ \uline{sur} $S$ qui a les propriétés
suivantes:
\[ \boxed{ f_{G/S} \circ \mathbb{V}_G = p \ \mathrm{Id}_{G^{(p)}} \qquad \mathbb{V}_G \circ \mathbb{f}_{G/S} = p \ \mathrm{Id}_G } \]
(Pour la définition de $\mathbb{V}_G$ et la démonstration, voir S.G.A. 3, VII$_A$
4.2-4.3).

Dans la plupart des cas, la seconde formule suffit pour calculer
$\mathbb{V}_G$. En effet, si $G$ est lisse alors $\mathbb{f}_{G/S}$ est un épimorphisme. Si $H$
se plonge dans un groupe lisse $G$, $\mathbb{V}_H$ est aussi déterminé car $\mathbb{V}_H$ est
induit par $\mathbb{V}_G$. La plupart des groupes que l'on considérera sont finis
localement libres sur la base, et on montrera plus loin qu'ils se plongent
canoniquement dans un groupe lisse.

Le morphisme $\mathbb{V}_G$ est fonctoriel en $G$ et commute aux changements
de base. Si $S = \mathrm{Spec}(\mathbb{F}_p)$ et si $G$ est le groupe $\mathbb{W}$ des vecteurs de
Witt, le morphisme de Frobenius et la Verschiebung coincident avec ceux
définis au §2.

\mspage{12}

\clearpage
% Notes of the Montreal seminar, 1970 (cote 129, typed part): archive cote 129, site page 13 (PDF page 14).
% typescript. Draft compared: none.
% AI-transcribed by Claude Opus 5.5 (claude-opus-5-5, adaptive thinking, effort high), Message Batches API, batch mtl129_full (msgbatch_01AGnDKHwy6ATCM7ehUa6263), 2026-11-13; not yet checked by eye.
\scanpage{129/13}
\margin{\unsure{$\cdot\ \cdot\ \nearrow$}}

5.4. \quad Pour tout $n \geq 1$, on définit les schémas
\[ X^{(p^n)} = X^{(p^n/S)} = [X^{(p^{n-1}/S)}]^{(p/S)} . \]
On note $\mathbb{F}^n_{X/S}$ le morphisme composé
\[ \mathbb{F}^n_{X/S} : X \xrightarrow{\mathbb{F}_{X/S}} X^{(p)} \xrightarrow{\mathbb{F}_{X^{(p)}/S}} X^{(p^2)} \to \dots \xrightarrow{\mathbb{F}_{X^{(p^{n-1})}/S}} X^{(p^n)} . \]
De même, si $G$ est un schéma en groupes commutatifs plat sur $S$, on
note $\mathbb{V}^n_G$ le composé
\[ \mathbb{V}^n_G : G^{(p^n)} \xrightarrow{\mathbb{V}_{G^{(p^{n-1})}}} G^{(p^{n-1})} \to \dots \xrightarrow{\mathbb{V}_G} G . \]
On a alors la conséquence facile des relations de 5.3 (correspondant
au cas $n = 1$):
\[ \mathbb{V}^n_G \circ \mathbb{F}^n_{G/S} = p^n \, Id_G \ , \quad \mathbb{F}^n_{G/S} \circ \mathbb{V}^n_G = p^n \, Id_{G^{(p^n)}} . \]
$\mathbb{F}^n_{X/S}$ et $\mathbb{V}^n_G$ s'appellent respectivement le \uline{$n^{\text{ième}}$ itéré du Frobenius}
\uline{relatif de} $G$ \uline{sur} $S$ et la \uline{$n^{\text{ième}}$ itéré de la Verschiebung de} $G$ \uline{sur} $S$.
Dans le cas où la base $S$ est le spectre du corps premier $\mathbb{F}_p$, ce sont
bien les itérés de $\mathbb{F}_{G/S}$, $\mathbb{V}_G$ au sens habituel.

\mspage{\hand{13}}

\clearpage
% Notes of the Montreal seminar, 1970 (cote 129, typed part): archive cote 129, site page 15 (PDF page 16).
% typescript. Draft compared: none.
% AI-transcribed by Claude Opus 5.5 (claude-opus-5-5, adaptive thinking, effort high), Message Batches API, batch mtl129_full (msgbatch_01AGnDKHwy6ATCM7ehUa6263), 2026-11-13; not yet checked by eye.
\scanpage{129/15}
CHAPITRE I

PRELIMINAIRES SUR WITT

\mspage{\hand{15}}

\clearpage
% Notes of the Montreal seminar, 1970 (cote 129, typed part): archive cote 129, site page 16 (PDF page 17).
% typescript. Draft compared: none.
% AI-transcribed by Claude Opus 5.5 (claude-opus-5-5, adaptive thinking, effort high), Message Batches API, batch mtl129_full (msgbatch_01AGnDKHwy6ATCM7ehUa6263), 2026-11-13; not yet checked by eye.
\scanpage{129/16}
\uline{PRELIMINAIRES SUR WITT}

1. \uline{Rappels sur les vecteurs de Witt}

Dans tout ce qui suit, $p$ désigne un nombre premier fixé une
fois pour toutes.

Le \uline{schéma des vecteurs de Witt} $\mathbb{W}$ et le \uline{schéma des vecteurs
de Witt tronqués à l'ordre $n$} $\mathbb{W}_n$ sont des schémas en anneaux, affines
sur $\mathbb{Z}$ qui sont définis comme suit :

Soit $n \geq 1$ un entier et soit $\mathbb{E}^n = \mathrm{Spec}\, \mathbb{Z}[T_1,T_2,\dots,T_n]$ d'es-%
pace affine de dimension $n$ sur $\mathbb{Z}$. Le schéma $\mathbb{E}^n$ représente le foncteur
$A \to A^n$ de la catégorie des schémas affines dans la catégorie des ensembles;
il est donc muni d'une structure naturelle de schéma en anneaux. On note
$\underline{O}^n$ le schéma $\mathbb{E}^n$ muni de cette structure. On démontre alors (c.f. J.-P.
Serre, ``Corps locaux'', II, §6) qu'il existe une unique structure de schéma
en anneaux sur $\mathbb{E}^n$ telle que le morphisme de schémas
\[ \phi = (\phi_1,\phi_2,\dots,\phi_n) : \mathbb{E}^n \to \underline{O}^n \]
défini par
\[ \phi_i(x_1,x_2,\dots,x_n) = x_1^{p^{i-1}} + p\, x_2^{p^{i-2}} + p^2\, x_3^{p^{i-3}} + \dots + p^{i-1}\, x_i \]
soit un homomorphisme d'anneaux.

Par définition, $\mathbb{W}_n$ est le schéma $\mathbb{E}_n$ muni de cette structure
de schéma en anneaux.

\mspage{16}

\clearpage
% Notes of the Montreal seminar, 1970 (cote 129, typed part): archive cote 129, site page 17 (PDF page 18).
% typescript. Draft compared: none.
% AI-transcribed by Claude Opus 5.5 (claude-opus-5-5, adaptive thinking, effort high), Message Batches API, batch mtl129_full (msgbatch_01AGnDKHwy6ATCM7ehUa6263), 2026-11-13; not yet checked by eye.
\scanpage{129/17}
Les morphismes de restriction
\[
R_n : \mathbb{W}_{n+1} \to \mathbb{W}_n \qquad (x_1,\ldots,x_n,x_{n+1}) \to (x_1,\ldots,x_n)
\]
sont alors des homomorphismes d'anneaux. En effet, il résulte de l'uni-%
cité de la structure d'anneau sur $\mathbb{W}_n$ qu'un morphisme $f : X \to \mathbb{W}_n$ est
un morphisme d'anneaux (resp. de groupes) si et seulement si $\phi_i \circ f : X \to \underline{O}$
est un homomorphisme d'anneaux (resp. de groupes pour tout $i$, $1 \leq i \leq n$ .
Ici, on a $\phi_i R_n = \phi_i$ , qui est par définition un homomorphisme d'anneaux.

On définit le schéma en anneau des vecteurs de Witt $\mathbb{W}$ par
\[
\mathbb{W} = \lim_{\substack{\leftarrow \\ n}} (\mathbb{W}_n , R_n) \quad .
\]
Le schéma sous-jacent à $\mathbb{W}$ est donc en particulier isomorphe à $\mathbb{E}^{\mathbb{N}}$ .

2. \uline{Quelques morphismes remarquables}

2.1. On définit
\[
V = V_n : \mathbb{W}_n \to \mathbb{W}_n \qquad (x_1,x_2,\ldots,x_{n-1},x_n) \to (0,x_1,\ldots,x_{n-1})
\]
\[
T = T_n : \mathbb{W}_n \to \mathbb{W}_{n+1} \qquad (x_1,x_2,\ldots,x_n) \to (0,x_1,\ldots,x_n) \ .
\]

Le morphisme $T$ est \uline{additif} ; on a en effet $\phi_i\, T = p\, \phi_{i-1}$
pour $1 \leq i \leq n+1$ .

On en déduit que $V_n = R_n\, T_n = T_{n-1}\, R_{n-1}$ est aussi un morphis-%
me \uline{additif}.

\mspage{\hand{17}}

\clearpage
% Notes of the Montreal seminar, 1970 (cote 129, typed part): archive cote 129, site page 18 (PDF page 19).
% typescript. Draft compared: none.
% AI-transcribed by Claude Opus 5.5 (claude-opus-5-5, adaptive thinking, effort high), Message Batches API, batch mtl129_full (msgbatch_01AGnDKHwy6ATCM7ehUa6263), 2026-11-13; not yet checked by eye.
\scanpage{129/18}
Par passage à la limite, les morphismes $V_n$ définissent un
morphisme additif
\[ V : \mathbb{W} \to \mathbb{W} \]
appelé \uline{Verschiebung} ou \uline{morphisme de décalage}.

2.2. \qquad Le \uline{morphisme de Frobenius} $F : \mathbb{W} \to \mathbb{W}$ est défini par passage
à la limite à partir des morphismes
\[ F_n : \mathbb{W}_n \to \mathbb{W}_n \qquad (x_1,x_2,\ldots,x_n) \to (x_1^p,x_2^p,\ldots,x_n^p) \ . \]

Ce morphisme n'est en général pas additif, mais par réduc-%
tion en caractéristique $p$ , il définit un \uline{homomorphisme d'anneau}
\[ F_{\mathbb{F}_p} : \mathbb{W}_{\mathbb{F}_p} \to \mathbb{W}_{\mathbb{F}_p} \qquad (\text{resp. } F_{n\,\mathbb{F}_p}) \]
où $\mathbb{F}_p$ désigne le corps premier $\mathbb{Z}/p\mathbb{Z}$ et $\mathbb{W}_{\mathbb{F}_p}$ désigne le schéma induit
sur $\mathbb{F}_p$ par changement de base. En effet, les deux lois de composition
sur $\mathbb{W}$ sont données par
\[ (x_1,x_2,\ldots,x_n,\ldots) + (y_1,\ldots,y_n,\ldots) = (S_1,S_2,\ldots,S_n,\ldots) \]
\[ (x_1,x_2,\ldots,x_n,\ldots) \cdot (y_1,\ldots,y_n,\ldots) = (P_1,\ldots,P_n,\ldots) \]
où les $S_i$ et $P_i$ sont des polynômes à coefficients entiers en
$(x_1,x_2,\ldots,x_i,\ y_1,y_2,\ldots,y_i)$ . Or, pour toute $\mathbb{F}_p$-algèbre, on a
\[ P(x^p,y^p) = [P(x,y)]^p \ . \]

\mspage{\hand{18}}

\clearpage
% Notes of the Montreal seminar, 1970 (cote 129, typed part): archive cote 129, site page 19 (PDF page 20).
% typescript. Draft compared: none.
% AI-transcribed by Claude Opus 5.5 (claude-opus-5-5, adaptive thinking, effort high), Message Batches API, batch mtl129_full (msgbatch_01AGnDKHwy6ATCM7ehUa6263), 2026-11-13; not yet checked by eye.
\scanpage{129/19}
\uline{Proposition 2.3.} En caractéristique $p$ , on a les relations
\[ F\cdot V = V\cdot F = p\ \mathrm{Id}_{\mathbb{W}_{\mathbb{F}_p}} \]
(c.f. J.-P. Serre ``Corps locaux'' II, §6, Cor. du th. 7). On peut donner
la démonstration directe suivante : Posons
\[ h = (h_1,h_2,\dots,h_n,\dots) = (p.\mathrm{Id}_W - VF)\ (x_1,\dots,x_n,\dots) . \]
On sait que $h_i \in \mathbb{Z}[x_1,x_2,\dots,x_i]$ , et il suffit de démontrer que $h_i$
est divisible par $p$ pour tout $i$ pour prouver la proposition. Or on
trouve que
\[ \phi_i(h) = p\phi_i(x) - \phi_i(VF(x)) = p^i\, x_i \]
d'où l'on déduit que $h_1 = px_1$ ,
\[ \phi_2(h) = h_1^p + ph_2 = p^2x_2 \ \text{ d'où }\ h_2 = p(x_2-p^{p-2}\, x_1) , \]
et par récurrence sur $i$ , on en déduit que pour tout $i$ on a $h_i = pk_i$ .

Notons en particulier que dans $\mathbb{W}_{\mathbb{F}_p}$ on a $p = (0,1,0,\dots)$
(car $(1,0,0,\dots)$ est l'unité de $\mathbb{W}$ ).

3. \uline{L'ind-schéma} $\underset{\rightarrow}{\mathbb{W}}$ et \uline{l'ind-schéma} $\underset{\rightarrow}{\mathbb{W}}'$

On définit
\[ \underset{\rightarrow}{\mathbb{W}} = \underset{\substack{\rightarrow\\ n}}{\lim}\ (\mathbb{W}_n,T_n) . \]
C'est un ind-schéma en groupes dont le ind-schéma sous-jacent est isomor-%

\mspage{19}

\clearpage
% Notes of the Montreal seminar, 1970 (cote 129, typed part): archive cote 129, site page 20 (PDF page 21).
% typescript. Draft compared: none.
% AI-transcribed by Claude Opus 5.5 (claude-opus-5-5, adaptive thinking, effort high), Message Batches API, batch mtl129_full (msgbatch_01AGnDKHwy6ATCM7ehUa6263), 2026-11-13; not yet checked by eye.
\scanpage{129/20}
phe à $\mathbb{E}^{(\mathbb{N})}$ , c'est-à-dire que pour tout anneau $A$ ,
\[
\underset{\to}{\mathbb{W}}(A) = A^{(\mathbb{N})} = \{(x_1,x_2,\ldots,x_n,\ldots) \ , \ x_i \in A \ , \ x_i = 0 \ \text{ sauf un nombre fini}\} \ .
\]
Mais ce foncteur a une structure additive qui n'est pas la structure
habituelle de $\mathbb{E}^{(\mathbb{N})}$ .

On peut définir un autre système inductif naturel à l'aide
des morphismes
\[
\mathbb{W}_n \xrightarrow{T_n F_n} \mathbb{W}_{n+1} \ .
\]

On a alors des diagrammes commutatifs
\begin{diagramblock}
\begin{tikzcd}
\mathbb{W}_n \arrow[r, "T_n"] \arrow[d, "F_n^{n-1}"] & \mathbb{W}_{n+1} \arrow[d, "F_{n+1}^{n}"] \\
\mathbb{W}_n \arrow[r, "T_n F_n = p"] & \mathbb{W}_{n+1}
\end{tikzcd}
\end{diagramblock} ,

car $T_n F_n = F_{n+1} T_n$ . Par définition, on pose
\[
\mathbb{W}' = \underset{\to}{\lim}(\mathbb{W}_n, T_n F_n) \ .
\]

En caractéristique $0$ , $\underset{\to}{\mathbb{W}}'$ n'a aucune structure naturelle de Groupe,
mais en caractéristique $p$ $\underset{\to}{\mathbb{W}}'$ est un ind-schéma en groupes, et on a un
morphisme additif
\[
u : \underset{\to}{\mathbb{W}}_{\mathbb{F}_p} \to \underset{\to}{\mathbb{W}}'_{\mathbb{F}_p}
\]
obtenu par passage à la limite des $F^{n-1}$ .

\mspage{\hand{20}}

\clearpage
% Notes of the Montreal seminar, 1970 (cote 129, typed part): archive cote 129, site page 21 (PDF page 22).
% typescript. Draft compared: none.
% AI-transcribed by Claude Opus 5.5 (claude-opus-5-5, adaptive thinking, effort high), Message Batches API, batch mtl129_full (msgbatch_01AGnDKHwy6ATCM7ehUa6263), 2026-11-13; not yet checked by eye.
\scanpage{129/21}
Notons qu'en caractéristique $p$ , les morphismes de transition $TF$ ne
sont autres que \uline{la multiplication par} $p$ puisque l'on a $TFR = FTR = FV = p.Id$ .

\uline{Remarque}. Si $A$ est un anneau \uline{parfait} de caractéristique $p$ , $u$ in-%
duit un isomorphisme
\[ u(A) : \mathbb{W}_{\underset{\to}{\mathbb{F}}_{p}}(A) \to \mathbb{W}'_{\underset{\to}{\mathbb{F}}_{p}}(A) \]
puisque c'est vrai pour chaque $F^{n}(A)$ .

4. \uline{L'anneau} $W$

4.1. On rappelle que (voir J.-P. Serre ``Corps locaux'' II, §6, th.7)
que si $k$ est un corps parfait de caractéristique $p$ , l'anneau
\[ W = \mathbb{W}(k) = \underset{\leftarrow}{\lim}\, W_{n} \qquad (Wn = \mathbb{W}_{n}(k)) \]
est un anneau de valuation discrète, complet, de corps résiduel $k$ , d'uni-%
formisante $p = (0,1,0,0,\ldots)$ , dont l'homomorphisme d'augmentation est
donné par la projection canonique
\[ W \to W_{1} = k . \]

On peut d'ailleurs montrer (loc. cit. II, §5, th. 3) que ces propriétés
caractérisent l'anneau $W$ à un isomorphisme unique près (théorème de Co-%
hen).

De plus, le corps des fractions de $W$ ,
\[ K = W[1/p] \]
est de caractéristique nulle.

\mspage{\hand{21}}

\clearpage
% Notes of the Montreal seminar, 1970 (cote 129, typed part): archive cote 129, site page 22 (PDF page 23).
% typescript. Draft compared: none.
% AI-transcribed by Claude Opus 5.5 (claude-opus-5-5, adaptive thinking, effort high), Message Batches API, batch mtl129_full (msgbatch_01AGnDKHwy6ATCM7ehUa6263), 2026-11-13; not yet checked by eye.
\scanpage{129/22}
Comme k est parfait, l'homomorphisme u définit un iso-%
morphisme de groupes
\[ u(k) : \underset{\rightarrow}{\mathbb{W}}(k) \rightarrow \underset{\rightarrow}{\mathbb{W}}'(k) \quad . \]

Remarquons alors que la projection naturelle $\pi_n : W \rightarrow W_n$
définit un isomorphisme
\[ \varepsilon_n : W/p^n\ W \overset{\sim}{\rightarrow} W_n \]
et l'on vérifie immédiatement que l'on a un diagramme commutatif
\begin{diagramblock}
\begin{tikzcd}
W \arrow[rrr, "p\ =\ VF"] \arrow[ddr, "\pi_n"'] \arrow[dr] & & & W \arrow[dl] \arrow[dd, "\pi_{n+1}"] \\
& W/p^n W \arrow[r, "p"] \arrow[dl, "\varepsilon_n", "\sim"'] & W/p^{n+1} W \arrow[dr, "\varepsilon_{n+1}"', "\sim"] & \\
W_n \arrow[rrr, "TF\ =\ FT"] & & & W_{n+1}
\end{tikzcd}
\end{diagramblock}

Utilisant alors l'isomorphisme naturel (de W-modules) valable pour tout
anneau de valuation discrète
\[ \underset{\rightarrow}{\lim}(W/p^n\ W,p) \overset{\sim}{\rightarrow} K/W \]
on déduit de ce qui précède que $\underset{\rightarrow}{\mathbb{W}}'(k)$ et donc $\underset{\rightarrow}{\mathbb{W}}(k)$ est isomorphe au
module dualisant $K/W$ . Comme $K/W$ est un W-module, on déduit une struc-%
ture de W-module sur $\underset{\rightarrow}{\mathbb{W}}(k)$ par transport de structure.

\uline{Proposition 4.2.} \unsure{Cetre|Cette} structure se prolonge de manière unique en une
structure de ind-schéma en groupes à anneau d'opérateurs $W$ sur
\[ \underset{\rightarrow}{W}_k = \underset{\rightarrow}{\lim}(\mathbb{W}_n)_k, T) \]

\mspage{\hand{22}}

\clearpage
% Notes of the Montreal seminar, 1970 (cote 129, typed part): archive cote 129, site page 23 (PDF page 24).
% typescript. Draft compared: none.
% AI-transcribed by Claude Opus 5.5 (claude-opus-5-5, adaptive thinking, effort high), Message Batches API, batch mtl129_full (msgbatch_01AGnDKHwy6ATCM7ehUa6263), 2026-11-13; not yet checked by eye.
\scanpage{129/23}
en faisant opérer $\lambda \in W$ sur $(\mathbb{W}_n)_k$ par multiplication par $F^{1-n}(\lambda)$ .

Remarquons que l'élément $F^{1-n}(\lambda) \in W$ est bien défini, car
$k$ étant parfait, $F : W \to W$ est un isomorphisme. Cette multiplication
a un sens car, pour toute k-algèbre $A$ , $\mathbb{W}_n(A)$ est une W-algèbre via le
morphisme
\[ W = \mathbb{W}(k) \xrightarrow{\pi_n} \mathbb{W}_n(k) \to \mathbb{W}_n(A) . \]

Pour que cette action se transporte à $\underset{\rightarrow}{\mathbb{W}}_k$ , il faut montrer
que l'action de $\lambda \in W$ commute aux morphismes $T$ . On vérifie immédia-%
tement qu'il suffit de montrer que pour toute k-algèbre $A$ , on a :
\[ \forall\ x \in \mathbb{W}_n(A) \ , \ \forall\ \lambda' \in \mathbb{W}_{n+1}(A) \ , \ \lambda'.T(x) = T(FR(\lambda').x) \]
ou, ce qui revient au même, puisque $R$ est une surjection :
\[ \forall\ \lambda' \ , \ y \in \mathbb{W}_{n+1}(A) \ , \ \lambda',V(y) = V(F(\lambda'),y) \ . \]
Or, si $A$ est un anneau parfait, $F$ est un isomorphisme, et cette rela-%
tion est vérifiée car
\[ F(\lambda',V(y)) = F(\lambda')\ FV(y) = F(\lambda').py = p[F(\lambda').y] = FV[F(\lambda').y] \ . \]
Dans le cas général, il suffit de remarquer que $\lambda'.V(y)$ et $V(F(\lambda').y)$
s'expriment à l'aide de polynômes à coefficients dans $\mathbb{Z}/p\mathbb{Z}$ . D'après ce
qui précède, ces polynômes prennent la même valeur si les variables sont
prises dans un anneau parfait, par exemple la clôture algébrique de $k$ .
On en déduit qu'ils sont égaux.

\mspage{\hand{23}}

\clearpage
% Notes of the Montreal seminar, 1970 (cote 129, typed part): archive cote 129, site page 24 (PDF page 25).
% typescript. Draft compared: none.
% AI-transcribed by Claude Opus 5.5 (claude-opus-5-5, adaptive thinking, effort high), Message Batches API, batch mtl129_full (msgbatch_01AGnDKHwy6ATCM7ehUa6263), 2026-11-13; not yet checked by eye.
\scanpage{129/24}
Enfin, il faut montrer que cette structure sur $\underset{\rightarrow}{\mathbb{W}}_k$ induit
sur $\underset{\rightarrow}{\mathbb{W}}(k)$ la structure précédemment définie. Or il est immédiat (et ce-%
la prouve l'unicité de la structure sur $\underset{\rightarrow}{\mathbb{W}}_k$) que la multiplication par
$F^{1-n}(\lambda)$ sur $W_n(k) = W_n$ induit sur $W/p^nW$ la multiplication par $\lambda$
via les isomorphismes
\[
\begin{diagramblock}
\begin{tikzcd}
W_n \arrow[r, "F^{n-1}", "\sim"'] & W_n & W/p^nW \arrow[l, "\varepsilon_n"', "\sim"] \ .
\end{tikzcd}
\end{diagramblock}
\]

\uline{Remarque}. Il n'y a pas en général d'action de $\mathbb{W}_k$ sur $\underset{\rightarrow}{\mathbb{W}}_k$ mais seu-%
lement de $W = \mathbb{W}(k)$ sur $\underset{\rightarrow}{\mathbb{W}}_k$ .

Cette construction est fonctorielle en $k$ et est variable
pour tout anneau parfait $A$ (mais $\mathbb{W}(A)$ n'est plus en général un an-%
neau de valuation discrète).

5. \uline{Les morphismes de Frobenius et de Verschiebung d'un schéma en groupes}

Dans tout ce qui suit, on suppose que tous les schémas sont
des schémas sur le corps premier $\mathbb{F}_p$ .

5.1. Si $X$ est un schéma, on appelle \uline{morphisme de Frobenius abso-}%
\uline{lu} de $X$ et on note
\[
\mathbb{F}_X : X \rightarrow X
\]
l'endomorphisme de $X$ qui est l'identité sur l'espace sous-jacent et qui
associe à une section $s$ de $\underline{O}_X$ la section $s^p$ . Si $Y$ est un autre
schéma, et $g : X \rightarrow Y$ un morphisme, le carré suivant est commutatif :

\mspage{24}

\clearpage
% Notes of the Montreal seminar, 1970 (cote 129, typed part): archive cote 129, site page 25 (PDF page 26).
% typescript. Draft compared: none.
% AI-transcribed by Claude Opus 5.5 (claude-opus-5-5, adaptive thinking, effort high), Message Batches API, batch mtl129_full (msgbatch_01AGnDKHwy6ATCM7ehUa6263), 2026-11-13; not yet checked by eye.
\scanpage{129/25}
\begin{diagramblock}\begin{tikzcd}
X \arrow[r, "\mathbb{F}_X"] \arrow[d, "g"'] & X \arrow[d, "g"'] \\
Y \arrow[r, "\mathbb{F}_Y"] & Y \quad .
\end{tikzcd}\end{diagramblock}

On obtient donc un endomorphisme $\mathbb{F}$ du foncteur identique de $\mathrm{Sch}/\mathbb{F}_p$ .

5.2. Soit $S$ un schéma fixé et $X$ un S-schéma. On note $X^{(p/S)}$
ou simplement $X^{(p)}$ le S-schéma image réciproque de $X$ par le change-%
ment de base $\mathbb{F}_S : S \to S$ .

Le diagramme ci-dessus étant commutatif, il existe une flè-%
che unique $\mathbb{F}_{X/S}$ telle que le diagramme suivant soit commutatif :

\begin{diagramblock}\begin{tikzcd}
X \arrow[dr] \arrow[drr, "\mathbb{F}_X"] \arrow[ddr, "\mathbb{F}_{X/S}"'] & & \\
 & X^{(p)} \arrow[r] \arrow[d] & X \arrow[d] \\
 & S \arrow[r, "\mathbb{F}_S"] & S \quad .
\end{tikzcd}\end{diagramblock}

On appellera le morphisme $\mathbb{F}_{X/S}$ ainsi défini le \uline{morphisme de Frobenius}
\uline{relatif de} $X$ \uline{au-dessus de} $S$ .

L'application $X \to \mathbb{F}_{X/S}$ est un morphisme fonctoriel du fonc-%
teur identique de $\mathrm{Sch}/S$ dans le foncteur $X \to X^{(p/S)}$ . Ce dernier fonc-%
teur commute aux limites projectives finies, donc si $X$ est un S-groupe,
$X^{(p/S)}$ l'est aussi et, $\mathbb{F}_{X/S}$ étant fonctoriel, est un morphisme de groupes.

La construction de $X^{(p/S)}$ commute aux changements de base,
c'est-à-dire que pour tout morphisme $T \to S$ , on a

\mspage{\hand{25}}

\clearpage
% Notes of the Montreal seminar, 1970 (cote 129, typed part): archive cote 129, site page 26 (PDF page 27).
% typescript. Draft compared: none.
% AI-transcribed by Claude Opus 5.5 (claude-opus-5-5, adaptive thinking, effort high), Message Batches API, batch mtl129_full (msgbatch_01AGnDKHwy6ATCM7ehUa6263), 2026-11-13; not yet checked by eye.
\scanpage{129/26}
\[ (X \underset{S}{\times} T)^{(p/T)} \cong X^{(p/S)} \underset{S}{\times} T \ . \]

Enfin, si $S$ est le spectre du corps premier $\mathbb{F}_p$ , $X^{(p/S)} = X$ et
$\mathbb{F}_{X/S} = \mathbb{F}_X$ .

5.3. \quad Soit $G$ un S-chéma en groupes commutatifs et plat variable.
On sait définir un homomorphisme canonique de groupes, fonctoriel en $G$
\[ \mathbb{V}_G : G^{(p/S)} \to G \]
appelé \uline{homomorphisme de Verschiebung de} $G$ \uline{sur} $S$ qui a les propriétés
suivantes :
\[ \boxed{ f_{G/S} \circ \mathbb{V}_G = p \ \mathrm{Id}_{G^{(p)}} \qquad \mathbb{V}_G \circ \mathbb{F}_{G/S} = p \ \mathrm{Id}_G } \]
(pour la définition de $\mathbb{V}_G$ et la démonstration, voir S.G.A. 3, VII$_A$,
4.2-4.3).

Dans la plupart des cas, la seconde formule suffit pour calcu-%
ler $\mathbb{V}_G$ . En effet, si $G$ est lisse alors $\mathbb{F}_{G/S}$ est un épimorphisme. Si
$H$ se plonge dans un groupe lisse $G$ , $\mathbb{V}_H$ est aussi déterminé car $\mathbb{V}_H$
est induit par $\mathbb{V}_G$ . La plupart des groupes que l'on considérera sont fi-%
nis localement libres sur la base, et on montrera plus loin qu'ils se plon-%
gent canoniquement dans un groupe lisse.

Le morphisme $\mathbb{V}_G$ est fonctoriel en $G$ et commute aux change-%
ments de base. Si $S = \mathrm{Spec}(\mathbb{F}_p)$ et si $G$ est le groupe $\mathbb{W}$ des vecteurs
de Witt, le morphisme de Frobenius et la Verschiebung coïncident avec ceux
définis au §2.

\mspage{26}

\clearpage
% Notes of the Montreal seminar, 1970 (cote 129, typed part): archive cote 129, site page 27 (PDF page 28).
% typescript. Draft compared: none.
% AI-transcribed by Claude Opus 5.5 (claude-opus-5-5, adaptive thinking, effort high), Message Batches API, batch mtl129_full (msgbatch_01AGnDKHwy6ATCM7ehUa6263), 2026-11-13; not yet checked by eye.
\scanpage{129/27}
5.4. \quad Pour tout $n \geq 1$ , on définit les schémas
\[ X^{(p^{n})} = X^{(p^{n}/S)} = [X^{(p^{n-1}/S)}]^{(p/S)} \quad . \]
On note $\mathbb{F}^{n}_{X/S}$ le morphisme composé
\[ \mathbb{F}^{n}_{X/S} : X \xrightarrow{\mathbb{F}_{X/S}} X^{(p)} \xrightarrow{\mathbb{F}_{X^{(p)}/S}} X^{(p^{2})} \to \dots \xrightarrow{\mathbb{F}_{X^{(p^{n-1})}/S}} X^{(p^{n})} \quad . \]
De même, si $G$ est un schéma en groupes commutatifs plat sur $S$ , on
note $\mathbb{V}^{n}_{G}$ le composé
\[ \mathbb{V}^{n}_{G} : G^{(p^{n})} \xrightarrow{\mathbb{V}_{G^{(p^{n-1})}}} G^{(p^{n-1})} \to \dots \xrightarrow{\mathbb{V}_{G}} G \ . \]
On a alors la conséquence facile des relations de 5.3. (correspondant
au cas $n = 1$) :
\[ \mathbb{V}^{n}_{G} \circ \mathbb{F}^{n}_{G/S} = p^{n}\,\mathrm{Id}_{G} \ , \ \mathbb{F}^{n}_{G/S} \circ \mathbb{V}^{n}_{G} = p^{n}\,\mathrm{Id}_{G^{(p^{n})}} \quad . \]
$\mathbb{F}^{n}_{X/S}$ et $\mathbb{V}^{n}_{G}$ s'appellent respectivement le \uline{$n^{\text{ième}}$ itéré du Frobenius re-}%
\uline{latif de} $G$ \uline{sur} $S$ et la \uline{$n^{\text{ième}}$ itéré de la Verschiebung de} $G$ \uline{sur} $S$ .
Dans le cas où la base $S$ est le spectre du corps premier $\mathbb{F}_{p}$ , ce sont
bien les itérés de $\mathbb{F}_{G/S}$ , $\mathbb{V}_{G}$ au sens habituel.

\mspage{\hand{27}}

\clearpage
% Notes of the Montreal seminar, 1970 (cote 129, typed part): archive cote 129, site page 28 (PDF page 29).
% typescript. Draft compared: none.
% AI-transcribed by Claude Opus 5.5 (claude-opus-5-5, adaptive thinking, effort high), Message Batches API, batch mtl129_full (msgbatch_01AGnDKHwy6ATCM7ehUa6263), 2026-11-13; not yet checked by eye.
\scanpage{129/28}
CHAPITRE II

GROUPES FINIS LOCALEMENT LIBRES ET THEORIE CLASSIQUE DE DIEUDONNE

\mspage{\hand{28}}

\clearpage
% Notes of the Montreal seminar, 1970 (cote 129, typed part): archive cote 129, site page 29 (PDF page 30).
% typescript. Draft compared: none.
% AI-transcribed by Claude Opus 5.5 (claude-opus-5-5, adaptive thinking, effort high), Message Batches API, batch mtl129_full (msgbatch_01AGnDKHwy6ATCM7ehUa6263), 2026-11-13; not yet checked by eye.
\scanpage{129/29}
\mspage{\hand{II.1}}

\uline{GROUPES FINIS LOCALEMENT LIBRES ET THEORIE CLASSIQUE DE DIEUDONNE}

1. \uline{Rappels généraux sur les groupes finis localement libres}

1.1. Soit $S$ un schéma et soit $G$ un schéma en groupes fini et
localement libre sur $S$ . Il est équivalent de dire que $G = \mathrm{Spec}\, A$ ,
où $A$ est une $\underline{O}_S$-algèbre quasi-cohérente finie localement libre, la
structure de groupe étant décrite par des homomorphismes de $\underline{O}_S$-algèbres
\[ \Delta : A \to A \otimes_{\underline{O}_S} A \]
\[ I : A \to A \]
\[ \varepsilon : A \to \underline{O}_S \]
où $\Delta$ , $I$ , $\varepsilon$ correspondent respectivement à la loi de composition, à
l'application $x \to x^{-1}$ et à la section unité, ces morphismes satisfai-%
sant des relations évidentes exprimant l'associativité, etc... Ainsi, $A$
muni de sa loi d'algèbre et de $\Delta$ devient une \uline{bigèbre}
% NOTE: a handwritten curved line runs from a caret between "algèbre" and "et" on the line above down around the following displayed formula, apparently meant to bring "mu : A (x) A -> A" in after "sa loi d'algèbre"
\margin{$\times$} \margin{$\checkmark$}
\[ \hand{\mu}\deleted{\illegible{1}} : A \otimes A \to A . \]

\margin{\hand{$\mathcal{J}$}} Le noyau \deleted{$J$}\inserted{$\mathcal{J}$} de $\varepsilon$ s'appelle \uline{l'idéal d'augmentation} du groupe $G$ .

\margin{\unsure{C.E.}} Dans tout ce qui suit, on supposera, sauf mention contraire,
que tous les groupes considérés sont \uline{commutatifs}.

1.2. Si $s \in S$ , on appelle \uline{rang de} $G$ \uline{en} $s$ le rang sur $k(s)$

\mspage{\hand{29}}

\clearpage
% Notes of the Montreal seminar, 1970 (cote 129, typed part): archive cote 129, site page 30 (PDF page 31).
% typescript. Draft compared: none.
% AI-transcribed by Claude Opus 5.5 (claude-opus-5-5, adaptive thinking, effort high), Message Batches API, batch mtl129_full (msgbatch_01AGnDKHwy6ATCM7ehUa6263), 2026-11-13; not yet checked by eye.
\scanpage{129/30}
\mspage{\hand{II.2}}

de $A \otimes_{\underline{O}_S} k(s)$ . L'application $s \to$ rang de $G$ en $s$ est une application
localement constante sur $S$ . Si ce rang est constant égal à
$n$ , on dit que $G$ est un \uline{groupe de rang} $n$ \uline{sur} $S$ .

1.3. \uline{Exemples de groupes finis localement libres}

1.3.1. Soit $G$ un groupe (abstrait) d'ordre $n$ . Soit $\underline{G}_S$ le fais-%
ceau constant sur $S$ défini par
\[ \underline{G}_S(S') = G \]
pour tout S-schéma $S'$ . $\underline{G}_S$ est un groupe libre de rang $n$ sur $S$ ; en
effet son algèbre affine est isomorphe à $\underline{O}_S^n$ .

1.3.2. Soit $\mathbb{G}_{m|S}$ le groupe multiplicatif sur $S$ , défini par
\[ \mathbb{G}_{m|S}(S') = \Gamma(S', \underline{O}_{S'})^* \]
pour tout S-schéma $S'$ . Pour tout $n \geq 1$ , on a l'isomorphisme
\[ n\ \mathrm{Id} : \mathbb{G}_{m|S} \to \mathbb{G}_{m|S} \qquad x \to x^n . \]
\margin{$\times$ \ $\mathbb{\mu}_{n|S}$}Par définition, $\deleted{\mu}_{n|S}$ , \uline{groupe des racines} $n^{\text{ièmes}}$ \uline{de l'unité sur} $S$ , est
le noyau de cet homomorphisme. C'est un groupe libre fini derang $n$ sur $S$ ;
en effet, son algèbre affine est isomorphe à $\underline{O}_S[T]/(T^n-1)$ .

1.3.3. Soit $\mathbb{G}_{a|S}$ le groupe additif sur $S$ , défini par
\[ \mathbb{G}_{a|S}(S') = \Gamma(S', \underline{O}_{S'}) \]
\margin{$\times$ \ $\mathbb{F}_p$}pour tout S-schéma $S'$ . Si $S$ est un schéma sur le corps premier $\mathbb{F}_p$ , on
a un homomorphisme de Frobenius

\margin{30}

\clearpage
% Notes of the Montreal seminar, 1970 (cote 129, typed part): archive cote 129, site page 31 (PDF page 32).
% typescript. Draft compared: none.
% AI-transcribed by Claude Opus 5.5 (claude-opus-5-5, adaptive thinking, effort high), Message Batches API, batch pilotnext (msgbatch_01A9sdu8SYq23bEtqX9tUncQ), 2026-11-13; not yet checked by eye.
\scanpage{129/31}
\mspage{\hand{II.3}}

\[ F : \mathbb{G}_{a|S} \to \mathbb{G}_{a|S} \qquad x \to x^{p} . \]

On désigne par $\alpha_{p|S}$ le noyau de $F$ . C'est un groupe libre fini de rang
$p$ sur $S$ , d'algèbre affine isomorphe à $\underline{\mathcal{O}}_{S}[T]/T^{p}$ .

1.4. \uline{Dualité de Cartier}

Pour tout schéma en groupes (commutatif) $G$ sur $S$ , on pose :
\[ G^{*} = \underline{\mathrm{Hom}}_{S\text{-}gr}(G,\mathbb{G}_{m|S}) . \]

\margin{$x \mapsto$} On définit ainsi un foncteur contravariant $G \over{\to}{\mapsto} G^{*}$ de la catégorie des
groupes sur $S$ dans elle-même ; ce foncteur commute aux changements de base.

Si $G = \mathrm{Spec}(A)$ est fini localement libre sur $S$ , on montre
(SGA 3, VII$_{A}$, 3.3.1) que l'on a
\[ G^{*} = \mathrm{Spec}(\overset{\hand{\vee}}{A}) \]
avec
\[ \overset{\hand{\vee}}{A} = \underline{\mathrm{Hom}}_{\underline{\mathcal{O}}_{S}\text{-mod}}(A,\underline{\mathcal{O}}_{S}) , \]
\margin{$\times$} la structure de \deleted{bi-algèbre}\inserted{bigèbre} sur $\overset{\hand{\vee}}{A}$ étant définie à l'aide des transposés des
morphismes $\Delta$ et $\mu$ de 1.1. Alors, $G^{*}$ est donc aussi un groupe fini lo-%
calement libre sur $S$ , et $G^{*}$ est de même rang que $G$ . De plus, l'homo-%
\margin{$\times$} morphisme fonctoriel canonique $G \to (G^{*})^{\hand{*}}$ déduit de la définition de $G^{*}$
est défini par l'isomorphisme $A \simeq (\overset{\hand{\vee}}{A})^{\vee}$ , et est donc un isomorphisme :
\[ G \overset{\sim}{\to} (G^{*})^{*} . \]

On définit donc ainsi une anti-équivalence de la catégorie des
$S$-schémas en groupes (commutatifs) finis localement libres avec elle-même,
appelée \uline{dualité de Cartier}.

\mspage{\hand{31}}

\clearpage
% Notes of the Montreal seminar, 1970 (cote 129, typed part): archive cote 129, site page 32 (PDF page 33).
% typescript. Draft compared: none.
% AI-transcribed by Claude Opus 5.5 (claude-opus-5-5, adaptive thinking, effort high), Message Batches API, batch mtl129_full (msgbatch_01AGnDKHwy6ATCM7ehUa6263), 2026-11-13; not yet checked by eye.
\scanpage{129/32}
\mspage{\hand{II.4}}

Si l'on suppose que  S  est un schéma sur  $\mathbb{F}_p$ , on a défini
(I.5) le schéma en groupes  $G^{(p/S)}$  et les morphismes de Frobenius rela-%
tif  $\mathbb{F}_{G/S}$  et de Verschiebung  $\mathbb{V}_G$ . Si  G  est fini localement libre sur
S , $G^{(p/S)}$  l'est aussi et on montre alors (SGA 3, VII$_A$, 4.3.3) que la
dualité de Cartier \uline{échange les morphismes de Frobenius relatif et de}
\uline{Verschiebung}, i.e. on a :
\[ (\mathbb{V}_G)^* = \mathbb{F}_{G^*/S} \quad ; \quad (\mathbb{F}_{G/S})^* = \mathbb{V}_{G^*} \quad . \]

1.5.   \uline{Exemples}

a) Soit  $\underline{G}_S$  le groupe fini sur  S  attaché au groupe fini
(abstrait)  G  (1.3.1). On vérifie aisément que l'algèbre affine de  $(\underline{G}_S)^*$
est isomorphe à  $\underline{O}_S[G]$ , la $\underline{O}_S$-algèbre du groupe  G , l'application dia-%
gonale provenant de l'application diagonale  $G \rightarrow G \times G$ . Si  $G = \mathbb{Z}/n\mathbb{Z}$ ,
on trouve
\[ (\mathbb{Z}/n\mathbb{Z})^*_S \simeq \mu_{n|S} \quad , \]
\margin{$\times$} qui résulte aussi \deleted{diversement} \inserted{directement} des définitions.

b) En caractéristique  p , on montre que l'on a
\[ (\alpha_{p|S})^* \simeq \alpha_{p|S} \quad . \]
\margin{$\times$ \quad $\alpha_{p|S}$}

1.6.   \uline{Quotients de groupes finis localement libres}

Soit  $\mu : G' \rightarrow G$  un monomorphisme de groupes finis localement
libres sur  S , et soit  $G'' = G/G'$  le faisceau quotient (pour la topolo-%
gie f.p.p.f.). On a donc une suite exacte de faisceaux
\margin{$\pi$}
\[ 0 \rightarrow G' \xrightarrow{\mu} G \xrightarrow{\deleted{p}\inserted{\pi}} G'' \rightarrow 0 \quad . \]

\margin{32}

\clearpage
% Notes of the Montreal seminar, 1970 (cote 129, typed part): archive cote 129, site page 33 (PDF page 34).
% typescript. Draft compared: none.
% AI-transcribed by Claude Opus 5.5 (claude-opus-5-5, adaptive thinking, effort high), Message Batches API, batch mtl129_full (msgbatch_01AGnDKHwy6ATCM7ehUa6263), 2026-11-13; not yet checked by eye.
\scanpage{129/33}
\mspage{\hand{II.5}}

On prouve (SGA 3, V, 4.1) que $G''$ est \uline{représentable par un}
\margin{$\pi$} \uline{schéma en groupes finis localement libre sur} $S$ , et que \deleted{$p$}\inserted{$\pi$} est fidèle-%
ment plat, de sorte que $G$ est fini localement libre sur $G''$ , de rang
relatif égal au rang de $G'$ . Par transitivité, on en conclut donc que
\[ \mathrm{rang}(G) = \mathrm{rang}(G')\cdot\mathrm{rang}(G'') . \]

Par ailleurs, on montre que si dans la suite exacte de fais-%
ceaux ci-dessus, on suppose $G'$ et $G''$ représentables par des schémas
en groupes finis localement libres sur $S$ , alors $G$ est lui aussi repré-%
sentable par un groupe fini localement libre sur $S$ . C'est-à-dire qu'une
extension de groupes finis localement libres est encore finie localement
libre.

1.7. Soit $\ell$ un nombre premier. On appelle \uline{composante $\ell$-primaire}
du groupe $G$ sur $S$ le sous-groupe
\[ C(\ell) = \varinjlim_{n} \mathrm{Ker}\{G \xrightarrow{\ell^{n}\mathrm{Id}} G\} \quad ; \]
\margin{$\times$ $\ell$} c'est le plus grand sous-groupe de \hand{$\ell$}-torsion de $G$ .

Si $G$ est fini localement libre sur $S$ , on prouve que $G$ est
localement annulé par un entier $n$ (SGA 3, VIII, 7.3), et donc chacune de
ses composantes $\ell$-primaires est localement annulée par une puissance de $\ell$
\margin{$\times$ $\ell$} (i.e. est un \hand{$\ell$}-groupe), et $G$ admet alors une décomposition canonique
(comme faisceau f.p.p.f.) en ses composantes $\ell$-primaires :
\[ G = \coprod_{\ell} G(\ell) . \]

\margin{33}

\clearpage
% Notes of the Montreal seminar, 1970 (cote 129, typed part): archive cote 129, site page 34 (PDF page 35).
% typescript. Draft compared: none.
% AI-transcribed by Claude Opus 5.5 (claude-opus-5-5, adaptive thinking, effort high), Message Batches API, batch mtl129_full (msgbatch_01AGnDKHwy6ATCM7ehUa6263), 2026-11-13; not yet checked by eye.
\scanpage{129/34}
\mspage{II.6}

Si $S$ est quasi-compact, $G$ est annulé par un entier $n$ ,
et cette décomposition est finie. Par ailleurs, les $G(\ell)$ sont des sché-%
mas en groupes finis localement libres, en tant que facteurs directs de
$G$ qui l'est.

Nous nous intéresserons ici aux groupes pour lesquels on a
$G = G(p)$ , c'est-à-dire aux groupes finis localement libres sur $S$ qui
sont localement annulés par une puissance de $p$ . Si $G$ est un tel p-%
groupe, son dual $G^*$ est aussi un p-groupe.

2. \uline{Types particuliers de groupes}

2.1. Soit $G$ un groupe fini localement libre sur une base quelcon-%
que $S$ .

2.1.1. On dit que $G$ est \uline{infinitésimal sur} $S$ (ou est \uline{radiciel}) si
l'une des conditions équivalentes suivantes est vérifiée :

- l'idéal d'augmentation $\mathcal{J}$ de $G$ (1.1) est localement nilpotent.

- $\forall\, s \in S$ , $G_s$ est connexe.

- $\forall\, s \in S$ , $G_s$ ne possède qu'un seul point.

- $\forall\, s \in S$ , $G_{\overline{s}}$ ne possède qu'un seul point, (où $G_{\overline{s}}$ désigne
la fibre géométrique de $G$ au-dessus de $s$).

2.1.2. On dit que $G$ est \uline{étale} au-dessus de $S$ si le morphisme struc-%
tural de $G$ est étale. Il est équivalent de dire que $\forall\, s \in S$ , $G_s$ est
étale sur $k(s)$ , ou encore que $G_{\overline{s}}$ est réduit pour tout $s \in S$ , ou en-%
\margin{X} core que $G$ est localement (pour la topologie étale, ou f.p.p.f., ou \hand{f.p.q.c.)}
isomorphe à un groupe constant.

\mspage{34}

\clearpage
% Notes of the Montreal seminar, 1970 (cote 129, typed part): archive cote 129, site page 35 (PDF page 36).
% typescript. Draft compared: none.
% AI-transcribed by Claude Opus 5.5 (claude-opus-5-5, adaptive thinking, effort high), Message Batches API, batch mtl129_full (msgbatch_01AGnDKHwy6ATCM7ehUa6263), 2026-11-13; not yet checked by eye.
\scanpage{129/35}
\mspage{\hand{II.7}}

2.1.3. \quad On dit que $G$ est \uline{unipotent} si $G^*$ est infinitésimal (voir
SGA 3, XVII, pour une définition générale des groupes unipotents).

2.1.4. \quad On dit que $G$ est de \uline{type multiplicatif} sur $S$ si $G^*$ est
étale. Cela équivaut au fait que pour tout $s \in S$ , $G_s$ est de type multi-%
plicatif, ou encore à ce que $G$ se plonge localement pour la topologie
étale dans $(\mathbb{G}_{m|S})^r$ .

\margin{X} 2.1.5. \quad On dit que $G$ est \uline{bi-infinitésimal} sur $S$ s'il est infini-%
tésimal et unipotent (i.e. $G$ et $G^*$ sont infinitésimaux).

2.2. \quad Si $G$ est étale et localement annulé par une puissance de $p$
sur un $F_p$-schéma $S$ , alors $G^*$ est infinitésimal. En d'autres termes,
sur une base de catactéristique $p$ , tout p-groupe de type multiplicatif
est infinitésimal. En effet, pour tout $s \in S$ , $G_{\overline{s}}$ est isomorphe au grou-%
pe constant $G(\overline{k})_{\overline{k}}$ (où l'on désigne par $\overline{k}$ la clôture algébrique du
corps résiduel en $s$). D'après 1.5.1 l'algèbre affine de $(G_{\overline{s}})^* = (G^*)_{\overline{s}}$
s'identifie à l'algèbre $\overline{k}[G(\overline{k})]$ du groupe $G(\overline{k})$ . On vérifie alors immé-%
diatement que l'idéal de $\overline{k}[G(\overline{k})]$ , formé des éléments $\sum\limits_{g\in G(\overline{k})} \lambda_g.g$ , $\lambda_g \in \overline{k}$
tels que $\sum\limits_{g} \lambda_g = 0$ , est maximal et nilpotent car $G(\overline{k})$ est annulé par une
puissance de $p$ et $\overline{k}$ est de caractéristique $p$ . Donc $(G^*)_{\overline{s}}$ se réduit
à un point et $G^*$ est infinitésimal.

2.3. \quad Soit
\[
0 \to G' \overset{\mu}{\to} G \overset{\pi}{\to} G'' \to 0
\]
une suite exacte de groupes finis localement libres. Alors $G$ est infini-%
tésimal (resp. étale, resp. unipotent, resp. de type multiplicatif, resp.

\margin{35}

\clearpage
% Notes of the Montreal seminar, 1970 (cote 129, typed part): archive cote 129, site page 36 (PDF page 37).
% typescript. Draft compared: none.
% AI-transcribed by Claude Opus 5.5 (claude-opus-5-5, adaptive thinking, effort high), Message Batches API, batch mtl129_full (msgbatch_01AGnDKHwy6ATCM7ehUa6263), 2026-11-13; not yet checked by eye.
\scanpage{129/36}
\mspage{\hand{II.8}}

bi-infinitésimal) si et seulement si $G'$ et $G''$ sont infinitésimaux
(resp. étales, resp. unipotents, resp. de type multiplicatif, resp. bi-%
infinitésimaux).

On vérifie aisément l'assertion dans le cas infinitésimal
et dans le cas étale en se ramenant au cas où $S$ est le spectre d'un
corps algébriquement clos. Les autres cas s'en déduisent par dualité de
Cartier.

2.4. Exemples de p-groupes en caractéristique $p$ .

$(\mathbb{Z}/p\mathbb{Z})_S$ est étale sur $S$ et $\mu_{p|S} = (\mathbb{Z}/p\mathbb{Z})^*_S$ est donc de
type multiplicatif. Par ailleurs, $\alpha_{p|S} = (\alpha_{p|S})^*$ est bi-infinitésimal.
Parmi ces 3 groupes, $\alpha_{p|S}$ et $(\mathbb{Z}/p\mathbb{Z})_S$ sont les groupes unipotents tan-%
disque les groupes infinitésimaux sont $\alpha_{p|S}$ et $\mu_{p|S}$ .

L'intérêt de ces trois groupes réside dans le résultat sui-%
vant : si $S$ est le spectre d'un corps algébriquement clos et de carac-%
téristique $p$ , tout groupe annulé par une puissance de $p$ et fini sur
$S$ admet une suite de composition finie
\margin{$\times$}
\[ G \supset G_0 \supset G_1 \supset G_2 \ldots \supset G_{n-1} \supset G_n \hand{$= 0$} \]
où les $G_i/G_{i+1}$ sont isomorphes à l'un des groupes $\alpha_p$ , $\mu_p$ , $\mathbb{Z}/p\mathbb{Z}$ . De
plus, il résulte de 2.4. que $G$ est infinitésimal (resp. étale, resp.
unipotent, resp. de type multiplicatif, resp. bi-infinitésimal) si et
seulement si tous les $G_i/G_{i+1}$ sont isomorphes à $\alpha_p$ ou $\mu_p$ (resp.
$\mathbb{Z}/p\mathbb{Z}$ , resp. $\mathbb{Z}/p\mathbb{Z}$ ou $\alpha_p$ , resp. $\mu_p$ , resp. $\alpha_p$ ). (cf. SGA 3, XVII, 1.7.
et 4.2.1).

\margin{36}

\clearpage
% Notes of the Montreal seminar, 1970 (cote 129, typed part): archive cote 129, site page 37 (PDF page 38).
% typescript. Draft compared: none.
% AI-transcribed by Claude Opus 5.5 (claude-opus-5-5, adaptive thinking, effort high), Message Batches API, batch mtl129_full (msgbatch_01AGnDKHwy6ATCM7ehUa6263), 2026-11-13; not yet checked by eye.
\scanpage{129/37}
\mspage{\hand{II.9}}

2.5. \quad Supposons que $S$ soit un $\mathbb{F}_p$-schéma. Pour tout groupe fini
localement libre sur $S$ on définit alors les morphismes de Frobenius
relatif et de Verschiebung (I,5). On a les critères suivants :

2.5.1. \quad $G$ est infinitésimal $<==>$ $\mathbb{F}_{G/S}$ est localement nilpotent.

2.5.2. \quad $G$ est unipotent $<==>$ $\mathbb{V}_G$ est localement nilpotent.

2.5.3. \quad $G$ est bi-infinitésimal $<==>$ $\mathbb{F}_{G/S}$ et $\mathbb{V}_G$ sont localement
nilpotents.

2.5.4. \quad $G$ est étale $<==>$ $\mathbb{F}_{G/S}$ est un isomorphisme.

2.5.5. \quad $G$ est de type multiplicatif $<==>$ $\mathbb{V}_G$ est un isomorphisme.

Il suffit de prouver les assertions 2.5.1. et 2.5.4., les
autres s'en déduisant, d'après 1.4. par dualité de cartier.

\uline{Démonstration de 2.5.1.}

Soit $\varepsilon : S \to G$ la section unité de $G$ ($G$ est considéré
comme schéma en ``ensembles pointés'' grâce à cette section $\varepsilon$ ). Nous al-%
\margin{X} lons caractériser le ``noyau'' de $\mathbb{F}^n_{G/S}$ pour tout S-schéma $G$ muni d'une
section. On vérifie immédiatement que le noyau de $\mathbb{F}^n_{G/S}$ est défini par
l'idéal engendré par $F^n(J) \subset \underline{O}_G$ , image de l'idéal d'augmentation
par le morphisme $x \to x^{p^n}$ ; désignons par $J_n$ cet idéal. On a évidem-%
ment $J_n \subset J^{p^n}$ . Si $J$ est de type fini, on peut supposer qu'il a moins
de $p^k$ générateurs, pour un entier $k$ convenable, et on vérifie immé-%
diatement que l'on a alors
\[ J^{p^{n-k}} \subset J_n , \]
de sorte que les $J_n$ forment un système cofinal aux $J^{p^n}$ .

\margin{37}

\clearpage
% Notes of the Montreal seminar, 1970 (cote 129, typed part): archive cote 129, site page 38 (PDF page 39).
% typescript. Draft compared: none.
% AI-transcribed by Claude Opus 5.5 (claude-opus-5-5, adaptive thinking, effort high), Message Batches API, batch mtl129_full (msgbatch_01AGnDKHwy6ATCM7ehUa6263), 2026-11-13; not yet checked by eye.
\scanpage{129/38}
\mspage{\hand{II.10}}

On en déduit donc que, si $G$ est localement de type fini sur
$S$ , \unsure{$\mathcal{F}_{G/S}$|$\hat{\mathcal{F}}_{G/S}$} est localement nilpotent si et seulement si $\mathcal{J}$ est localement
nilpotent. Pour un groupe fini localement libre sur $S$ cela équivaut au
fait qu'il est infinitésimal.

\uline{Démonstration de 2.5.4.}

On montre d'une manière générale (SGA 5, XIV (Houzel)) que si
$G$ est un schéma localement de présentation finie sur $S$ , \unsure{$\mathcal{F}_{G/S}$|$\hat{\mathcal{F}}_{G/S}$} est un
isomorphisme si et seulement si $G$ est étale sur $S$ .

3. \uline{Décomposition d'un p-groupe fini localement libre sur une base de carac-}
\uline{téristique} $p$ \uline{réduite à un point} $s$ .

3.1. $G$ étant un groupe fini localement libre sur $S$ , on a une sui-%
te exacte canonique, fonctorielle en $G$
\[ (*) \qquad 0 \rightarrow G^{0} \rightarrow G \rightarrow G^{et} \rightarrow 0 \]
où $G^{0}$ est la composante neutre de $G$ et $G^{et} = G/G^{0}$ . Le sous-groupe $G^{0}$
est ouvert dans $G$ (SGA$_{3}$, VI$_{B}$, 3.9), il est donc fini localement libre sur
$S$ et c'est le plus grand sous-groupe infinitésimal de $G$ . Le groupe quo-%
tient $G^{et}$ est étale, (SGA 3, VI$_{A}$, 5.5) et il est fini localement libre
sur $S$ , d'après 1.6.

Appliquant la même construction à $G^{*}$ , on obtient un sous-grou-%
pe $G^{mult} = (G^{*}/(G^{*})^{0})^{*}$ , qui est le plus grand sous-groupe de type multi-%
plicatif de $G$ .

\mspage{\hand{38}}

\clearpage
% Notes of the Montreal seminar, 1970 (cote 129, typed part): archive cote 129, site page 39 (PDF page 40).
% typescript. Draft compared: none.
% AI-transcribed by Claude Opus 5.5 (claude-opus-5-5, adaptive thinking, effort high), Message Batches API, batch mtl129_full (msgbatch_01AGnDKHwy6ATCM7ehUa6263), 2026-11-13; not yet checked by eye.
\scanpage{129/39}
\mspage{\hand{II.11}}

Si $k(s)$ est de caractéristique $p$ et si $G$ est annulé
par une puissance de $p$ , $G^{\mathrm{mult}}$ est infinitésimal (2.2) et on a donc
des inclusions
\[ 0 \subset G^{\mathrm{mult}} \subset G^{0} \subset G \]
\margin{X} où $G/G^{0} = G^{\mathrm{et}}$ est étale, $G^{0}/G^{\mathrm{mult}}$ est bi-infinitésimal\hand{,}\deleted{(}car c'est
un quotient d'un groupe infinitésimal et son dual s'identifie à $((G^{0})^{*})^{0}$ .

Cette suite de composition de $G$ est fonctorielle en $G$ ,
et on vérifie qu'elle commute aux produits finis.

3.2. Si $S$ est le \uline{spectre d'un corps parfait de caractéristique} $p$ ,
la projection canonique $G \to G^{\mathrm{et}}$ induit un isomorphisme de $G_{\mathrm{red}}$ sur
$G^{\mathrm{et}}$ et la suite exacte $(*)$ splitte canoniquement (SGA 3, XVII, 1.6),
et on a donc
\[ G = G^{0} \times G^{\mathrm{et}} . \]

De même la suite exacte
\[ 0 \to G^{\mathrm{mult}} \to G^{0} \to G^{0}/G^{\mathrm{mult}} \to 0 \]
est splittée car la suite duale l'est. On obtient donc une décomposition
canonique, fonctorielle en $G$ :
\[ G = G^{\mathrm{mult}} \times G^{\mathrm{bi}} \times G^{\mathrm{et}} \]
où $G^{\mathrm{mult}}$, $G^{\mathrm{bi}} = G^{0}/G^{\mathrm{mult}}$ , $G^{\mathrm{et}}$ sont respectivement de type multiplica-%
tif, bi-infinitésimal et étale. La partie unipotente de $G$ est alors
$G^{\mathrm{bi}} \times G^{\mathrm{et}}$ et la partie infinitésimale est $G^{\mathrm{mult}} \times G^{\mathrm{bi}}$ .

\mspage{\hand{39}}

\clearpage
% Notes of the Montreal seminar, 1970 (cote 129, typed part): archive cote 129, site page 40 (PDF page 41).
% typescript. Draft compared: none.
% AI-transcribed by Claude Opus 5.5 (claude-opus-5-5, adaptive thinking, effort high), Message Batches API, batch mtl129_full (msgbatch_01AGnDKHwy6ATCM7ehUa6263), 2026-11-13; not yet checked by eye.
\scanpage{129/40}
\mspage{\hand{II.12}}

On a les équivalences :

$G$ est infinitésimal $\Longleftrightarrow G^{\mathrm{et}} = 0$

$G$ est étale $\Longleftrightarrow G^{\mathrm{mult}} = G^{\mathrm{bi}} = 0$

$G$ est unipotent $\Longleftrightarrow G^{\mathrm{mult}} = 0$

$G$ est bi-infinitésimal $\Longleftrightarrow G^{\mathrm{mult}} = G^{\mathrm{et}} = 0$

$G$ est de type multiplicatif $\Longleftrightarrow G^{\mathrm{bi}} = G^{\mathrm{et}} = 0$ .

3.3. La catégorie des p-groupes étales finis sur un corps parfait
$k$ est bien connue. Elle est équivalente à la catégorie des p-groupes fi-%
nis sur lesquels opère le groupe fondamental $\pi_1(k) = \mathrm{Gal}(\overline{k}/k)$ . Par dua-%
lité de Cartier, la catégorie des p-groupes finis de type multiplicatif
est équivalente à la catégorie opposée de celle des p-groupes étales. Par
ailleurs, la catégorie des p-groupes finis sur lesquels opère $\pi_1(k)$ est
équivalente à sa propre opposée par dualité de Prontryagin. On en déduit
une équivalence de catégories entre la catégorie des p-groupes finis de
type multiplicatif avec celle des p-groupes finis sur lesquels opère $\pi_1(k)$ .

4. \uline{Théorie de Dieudonné sur un corps parfait}

4.1. Soit $k$ un corps parfait de caractéristique $p > 0$ et soit
p-Gr-fin(k) la catégorie des schémas en groupes finis sur $k$ et annulés
par une puissance de $p$ . Le but de la théorie de Dieudonné est d'établir
une anti-équivalence entre la catégorie p-Gr-fin(k) et une catégorie de
modules que nous allons définir.

On appelle \uline{anneau de Dieudonné} \hand{$\mathcal{D}$} l'anneau $W[F,V]$ (non com-%
mutatif) engendré sur l'anneau local $W = W(k)$ (I,4) par des indéterminées

\margin{40}

\clearpage
% Notes of the Montreal seminar, 1970 (cote 129, typed part): archive cote 129, site page 41 (PDF page 42).
% typescript. Draft compared: none.
% AI-transcribed by Claude Opus 5.5 (claude-opus-5-5, adaptive thinking, effort high), Message Batches API, batch mtl129_full (msgbatch_01AGnDKHwy6ATCM7ehUa6263), 2026-11-13; not yet checked by eye.
\scanpage{129/41}
\mspage{\hand{II.13}}

F et V satisfaisant les relations suivantes :
\[ F\cdot\lambda = \lambda^{\sigma}\cdot F \]
\[ \lambda\cdot V = V\cdot\lambda^{\sigma} \]
\[ F\cdot V = V\cdot F = p \ , \]
pour tout $\lambda \in W$ et où l'on note par $\lambda^{\sigma}$ l'image de $\lambda$ par le morphis-%
me de Frobénius de $W$ .

On appelle \uline{catégorie des modules de Dieudonné} la catégorie
$\underline{\mathrm{Mod}}_{\hand{\mathcal{D}};W\ \hand{\ell.f.}}$ des \hand{$\mathcal{D}$}-modules à gauche tels que le W-module sous-jacent soit
de longueur finie. On voit en particulier qu'un tel module est annulé par
une puissance de $p$ .

Un objet de $\underline{\mathrm{Mod}}_{\hand{\mathcal{D}};W\ \hand{\ell.f.}}$ est donc un W-module $M$ de longueur
finie, muni de deux morphismes de W-modules
\[ F_M : M^{\sigma} \rightarrow M \]
\[ V_M : M \rightarrow M^{\sigma} \quad , \]
en posant
\margin{$\times$ \quad $\otimes$}
\[ M^{\sigma} = M \hand{\otimes}_W (W,\sigma) \ , \]
($\sigma$ étant le morphisme de Frobénius de $W$), $F_M$ et $V_M$ vérifiant
\[ F_M\cdot V_M = p\cdot \mathrm{Id}_M \ , \ V_M F_M = p\cdot \mathrm{Id}_M \quad . \]

\uline{Théorème\unsure{.}4.2.} (Dieudonné)

Il existe une équivalence de catégories
\[ D^* : [p\text{-}\mathrm{Gr}\text{-}\mathrm{fin}(k)]^{0} \hand{\xrightarrow{\ \approx\ }} \underline{\mathrm{Mod}}_{\hand{\mathcal{D}},W\ \hand{\ell.f.}} \]
telle que l'on ait un isomorphisme canonique $[D^*(G)]^{\unsure{0|\sigma}} \hand{\approx} D^*(G^{(p)})$ ,

\margin{41}

\clearpage
% Notes of the Montreal seminar, 1970 (cote 129, typed part): archive cote 129, site page 42 (PDF page 43).
% typescript. Draft compared: none.
% AI-transcribed by Claude Opus 5.5 (claude-opus-5-5, adaptive thinking, effort high), Message Batches API, batch mtl129_full (msgbatch_01AGnDKHwy6ATCM7ehUa6263), 2026-11-13; not yet checked by eye.
\scanpage{129/42}
\mspage{II.14}

\margin{$\times$} \hand{conduisant aux identifications}
\[ D^*(\over{F}{\mathbb{F}}_{G/k}) = F_{D^*(G)} \]
\[ D^*(\over{V}{\mathbb{V}}_G) = V_{D^*(G)} \ . \]

\uline{Corollaire 4.2.1}. On a les équivalences suivantes :

$G$ est infinitésimal $\Longleftrightarrow$ $F$ opère de façon nilpotente sur $D^*(G)$ .

$G$ est unipotent $\Longleftrightarrow$ $V$ opère de façon nilpotente sur $D^*(G)$ .

$G$ est bi-infinitésimal $\Longleftrightarrow$ $F$ et $V$ opèrent de façon nilpotente.

$G$ est étale $\Longleftrightarrow$ $F$ opère bijectivement sur $D^*(G)$ .

$G$ est de type multiplicatif $\Longleftrightarrow$ $V$ opère bijectivement sur $D^*(G)$ .

Pour une démonstration complète du théorème, se reporter
au livre de P. Gabriel et M. Demazure, ``Groupes algébriques''. Dans ce
qui suit, on donnera seulement la construction du foncteur $D^*$ .

On observe tout d'abord que, d'après 3.2., la catégorie
p\_Gr-fin(k) est équivalente à la catégorie produit de la sous-catégo-%
rie pleine de p-Gr-fin(k) dont les objets sont les groupes unipotents
et de la sous-catégorie pleine de p-Gr-fin(k) dont les objets sont
les groupes de type multiplicatif. Il suffira donc de construire $D^*$
séparément sur ces deux sous-catégories.

4.3. \uline{Cas des groupes unipotents}

4.3.1. Pour tout groupe unipotent sur $k$ , on pose par définition
\[ D^*(G) = \mathrm{Hom}_{k\text{-}Gr}(G,\underset{\rightarrow}{\mathbb{W}}) \]
où, pour abréger, on écrit $\underset{\rightarrow}{\mathbb{W}}$ pour $\underset{\rightarrow}{\mathbb{W}}_k$ (I,3). On sait que $W$ agit sur
$\underset{\rightarrow}{\mathbb{W}}$ (I,4.2) et donc $D^*(G)$ est un $W$-module, qui est de longueur finie si
$G$ est fini sur $k$ .

\mspage{42}

\clearpage
% Notes of the Montreal seminar, 1970 (cote 129, typed part): archive cote 129, site page 43 (PDF page 44).
% typescript. Draft compared: none.
% AI-transcribed by Claude Opus 5.5 (claude-opus-5-5, adaptive thinking, effort high), Message Batches API, batch mtl129_full (msgbatch_01AGnDKHwy6ATCM7ehUa6263), 2026-11-13; not yet checked by eye.
\scanpage{129/43}
\mspage{II.15}

L'action de $F$ et $V$ sur $D^*(G)$ se définit en posant,
pour tout $u \in D^*(G)$
\[ F\cdot u = F_{\underset{\rightarrow}{W}} \circ u \]
\[ V\cdot u = V_{\underset{\rightarrow}{W}} \circ u \]
où $F_{\underset{\rightarrow}{W}}$ et $V_{\underset{\rightarrow}{W}}$ sont les endomorphismes de Frobénius et de Verschiebung
\margin{XX} de $\underset{\rightarrow}{\mathbb{W}}$ , définis par passage à la limite des morphismes correspondant\hand{s} sur \inserted{les} $W_n$ .

On vérifie immédiatement, en revenant à la définition de l'ac-%
tion de $W$ sur $\underset{\rightarrow}{W}$ et en utilisant les propriétés de $F_{\underset{\rightarrow}{W}}$ et $V_{\underset{\rightarrow}{W}}$ , que
ces actions font de $D^*(G)$ un $\mathcal{D}$-module.

4.3.2. Montrons que l'on a un isomorphisme canonique
\[ D^*(G^{(p)}) \overset{\sim}{\rightarrow} (D^*(G))^{\sigma} \]
tel que
\[ D^*(\mathbb{F}_{G/k}) = F_{D^*(G)} \]
\[ D^*(\mathbb{V}_{G}) = V_{D^*(G)} \quad . \]

En effet, $k$ étant parfait, le changement de base $k \overset{\sigma}{\rightarrow} k$
induit un isomorphisme $\sigma$-linéaire
\[ \mathrm{Hom}(G, \mathbb{W}) \overset{\sim}{\rightarrow} \mathrm{Hom}(G^{(p)}, \mathbb{W}^{(p)}) \quad . \]

Mais $\underset{\rightarrow}{\mathbb{W}}$ étant défini sur le corps premier $\mathbb{F}_p$ , on a $\underset{\rightarrow}{\mathbb{W}}^{(p)} = \underset{\rightarrow}{\mathbb{W}}$
(car la formation de $\underset{\rightarrow}{\mathbb{W}}^{(p)}$ commute aux changements de base, cf. I, 5.2).
Par ailleurs
\[ D^*(\mathbb{F}_{G/k}) = \mathrm{Hom}(G^{(p)}, \underset{\rightarrow}{\mathbb{W}}) \rightarrow \mathrm{Hom}(G, \underset{\rightarrow}{\mathbb{W}}) \]
\margin{$\times$ \quad $\mapsto$}
\[ u \mapsto u \circ \mathbb{F}_{G/k} \]

\margin{43}

\clearpage
% Notes of the Montreal seminar, 1970 (cote 129, typed part): archive cote 129, site page 44 (PDF page 45).
% typescript. Draft compared: none.
% AI-transcribed by Claude Opus 5.5 (claude-opus-5-5, adaptive thinking, effort high), Message Batches API, batch mtl129_full (msgbatch_01AGnDKHwy6ATCM7ehUa6263), 2026-11-13; not yet checked by eye.
\scanpage{129/44}
\mspage{\hand{II.16}}

s'identifie à $F_{D^*(G)}$ car $u \circ \mathbb{F}_{G/k} = \mathbb{F}_{\underset{\rightarrow}{W}} \circ u$ et $\mathbb{F}_{\underset{\rightarrow}{W}}$ n'est autre que
$F_{\underset{\rightarrow}{W}}$ (I, 5.3).

On fait un raisonnement analogue pour établir l'assertion
relative à $D^*(\mathbb{V}_G)$ .

Si $G$ est unipotent, il résulte de ce que $\mathbb{V}_G$ est nilpotent
(2.5.2) que $V_{D^*(G)}$ opère de façon nilpotente. On peut alors montrer,
plus généralement, que le foncteur $D^*$ définit une équivalence de ca-%
tégories entre la catégorie des p-groupes affines unipotents sur $k$ avec
la catégorie des \hand{$\mathcal{D}$}-modules sur lesquels $V$ est localement nilpotent. On
a en particulier :

\uline{Corollaire 4.3.3} On a des équivalences de catégories :

\begin{diagramblock}\[
\left(\begin{array}{l} \text{p-groupes finis} \\ \text{unipotents sur } k \end{array}\right)^0 , \xrightarrow[\hand{\approx}]{D^*} \left(\begin{array}{l} \text{sous-catégorie pleine de } \underline{\mathrm{Mod}}_{\hand{\mathcal{D}},W\,\hand{\ell.f}} \\ \text{formée des objets } M \text{ où } V \text{ est nilpotent} \end{array}\right)
\]\end{diagramblock}

\begin{diagramblock}\[
\left(\begin{array}{l} \text{p-groupes finis} \\ \text{étales sur } k \end{array}\right)^0 \xrightarrow[\hand{\approx}]{D^*} \left(\begin{array}{l} \text{sous-catégorie pleine de } \underline{\mathrm{Mod}}_{\hand{\mathcal{D}},W\,\hand{\ell.f.}} \\ \text{formée des objets } M \text{ où } F \text{ est un isomorphisme} \end{array}\right.
\]\end{diagramblock}

\begin{diagramblock}\[
\left(\begin{array}{l} \text{p-groupes finis} \\ \text{bi-infinitésimaux} \end{array}\right)^0 \xrightarrow[\hand{\approx}]{D^*} \left(\begin{array}{l} \text{sous-catégorie pleine de } \underline{\mathrm{Mod}}_{\hand{\mathcal{D}},W\,\hand{\ell.f.}} \\ \text{formée des objets où } F \text{ et } V \text{ sont nilpotents} \end{array}\right)
\]\end{diagramblock}

Cela résulte de la description de $F$ et $V$ (4.3.2) et de
la caractérisation de chaque type de groupes à l'aide des morphismes $\mathbb{F}$
et $\mathbb{V}$ . (2.5).

\margin{44}

\clearpage
% Notes of the Montreal seminar, 1970 (cote 129, typed part): archive cote 129, site page 45 (PDF page 46).
% typescript. Draft compared: none.
% AI-transcribed by Claude Opus 5.5 (claude-opus-5-5, adaptive thinking, effort high), Message Batches API, batch mtl129_full (msgbatch_01AGnDKHwy6ATCM7ehUa6263), 2026-11-13; not yet checked by eye.
\scanpage{129/45}
\mspage{\hand{II.17}}

4.4. \uline{Cas des groupes de type multiplicatif}

On désire construire une équivalence de catégories.

\begin{diagramblock}\[
\left(\begin{array}{l}
\text{p-groupes finis sur } k\\
\text{de type multiplicatif}
\end{array}\right)^{0}
\xrightarrow[\hspace{2em}]{D^{*}\ \hand{$\approx$}}
\left(\begin{array}{l}
\text{sous catégorie pleine de } \mathrm{Mod}_{\hand{$\mathcal{D}$},W\ \hand{$\ell.f.$}}\\
\text{formée des objets où } V \text{ est un}\\
\text{isomorphisme}
\end{array}\right)
\]\end{diagramblock}

La catégorie de droite est équivalente à la catégorie des
W-modules de longueur finie munis d'un isomorphisme (W-linéaire)
\[ M \xrightarrow{\ V\ } M^{\sigma} = M \otimes_{W} (W,\sigma) \]
(en effet, F est alors déterminé de façon unique par la relation FV = p).
En considérant l'isomorphisme inverse, cette catégorie est encore équiva-%
lente à la catégorie des W-modules de longueur finie munis d'un isomorphis-%
me
\[ M^{\sigma} \xrightarrow{\ V^{-1}\ } M \ . \]

Mais, d'après le corollaire 4.3.3, cette catégorie est anti-%
équivalente à la catégorie des p-groupes finis étales sur k , qui est elle-%
même anti-équivalente, par dualité de Cartier, à la catégorie des p-groupes
finis de type multiplicatif sur k .

Pour obtenir en définitive la bonne variance pour le foncteur
D* cherché, il suffit de remarquer qu'il existe une anti-équivalence de
la catégorie des p-groupes finis étales sur k avec elle-même, donnée par
la dualité de Pontryagin
\[ G \to \underline{\mathrm{Hom}}_{\mathrm{gr}}(G,\ \mathbb{Q}_{p}/\mathbb{Z}_{p}) \ . \]
\margin{x $\mathbb{Q}_p/\mathbb{Z}_p$}

\mspage{\hand{45}}

\clearpage
% Notes of the Montreal seminar, 1970 (cote 129, typed part): archive cote 129, site page 46 (PDF page 47).
% typescript. Draft compared: none.
% AI-transcribed by Claude Opus 5.5 (claude-opus-5-5, adaptive thinking, effort high), Message Batches API, batch pilotnext (msgbatch_01A9sdu8SYq23bEtqX9tUncQ), 2026-11-13; not yet checked by eye.
\scanpage{129/46}
\mspage{II.18}

Le foncteur $D^*$ recherché est alors donné par la composition
de toutes ses équivalences. La situation est résumée par le diagramme
suivant :

\begin{diagramblock}
\begin{tikzcd}[column sep=large, row sep=large]
\left(\begin{array}{l}\text{p-groupes finis de}\\ \text{type multiplicatif}\end{array}\right)^{0}
\arrow[r, dashed, "D^*", "\hand{$\approx$}"']
\arrow[d, "\hand{$\wr$}"', "\text{dualité de Cartier}"]
&
\left(\begin{array}{l}\text{sous catégorie pleine de } \underline{\mathrm{Mod}}_{\hand{$\mathcal{D}$};\ W\ \hand{$\ell.f.$}}\\ \text{formée des objets où } V \text{ est } \deleted{objectif}\ \inserted{bijectif}\end{array}\right)
\arrow[d, "\hand{$\wr$}"]
\\
\left(\begin{array}{l}\text{p-groupes finis}\\ \text{étales}\end{array}\right)
\arrow[d, "\hand{$\wr$}"', "\text{dualité de Pontryagin}"]
&
\left(\begin{array}{l}\text{W-modules de longueur finie munis d'un}\\ \text{isomorphisme } V : M \to M^{\sigma}\end{array}\right)
\arrow[d, "\hand{$\wr$}"]
\\
\left(\begin{array}{l}\text{p-groupes finis}\\ \text{étales}\end{array}\right)^{0}
\arrow[r, "D^*", "\hand{$\sim$}"']
&
\left(\begin{array}{l}\text{sous catégorie pleine de } \underline{\mathrm{Mod}}_{\hand{$\mathcal{D}$};\ W\ \hand{$\ell.f.$}}\\ \text{formée des objets où } F \text{ est bijectif}\end{array}\right)
\end{tikzcd}
\end{diagramblock}
\margin{X}

(4.3.3)

4.5. \uline{Cas général}

Si $G$ est un p-groupe fini quelconque, on écrit,(d'après 3.2)
\[ G = G^{\mathrm{uni}} \times G^{\mathrm{mult}} \]
où $G^{\mathrm{uni}}$ et $G^{\mathrm{mult}}$ sont respectivement unipotent et de type multiplica-%
tif, et on pose
\[ D^*(G) = D^*(G^{\mathrm{uni}}) \times D^*(G^{\mathrm{mult}}) \ . \]

Le théorème de Dieudonné 4.2. sera démontré si l'on prouve que
l'on a une équivalence de catégories

\mspage{46}

\clearpage
% Notes of the Montreal seminar, 1970 (cote 129, typed part): archive cote 129, site page 47 (PDF page 48).
% typescript. Draft compared: none.
% AI-transcribed by Claude Opus 5.5 (claude-opus-5-5, adaptive thinking, effort high), Message Batches API, batch mtl129_full (msgbatch_01AGnDKHwy6ATCM7ehUa6263), 2026-11-13; not yet checked by eye.
\scanpage{129/47}
\mspage{\hand{II 19}}

\margin{$\times$}
\[
\underline{\mathrm{Mod}}_{\hand{$\mathcal{D}$};W\,\hand{$\ell.f.$}} \xrightarrow{\ \approx\ }
\left(\begin{array}{l}
\text{sous-catégorie pleine}\\
\text{de } \underline{\mathrm{Mod}}_{\hand{$\mathcal{D}$};W\,\hand{$\ell.f.$}} \text{ formée}\\
\text{des objets où } V \text{ est}\\
\text{nilpotent}
\end{array}\right)
\mathrm{x}
\left(\begin{array}{l}
\text{sous catégorie pleine}\\
\text{de } \underline{\mathrm{Mod}}_{\hand{$\mathcal{D}$};W\,\hand{$\ell.f.$}} \text{ formée}\\
\text{des objets où } V \text{ est}\\
\text{bijectif}
\end{array}\right) .
\]

Cela résulte du lemme suivant :

\uline{Lemme}. Tout \hand{$\mathcal{D}$}-module $M$ de longueur finie sur $W$ se décompose d'une
manière unique (et fonctorielle) en une somme directe
\[
M = M' \;\mathrm{x}\; M''
\]
telle que $V|M'$ est nilpotent et $V|M''$ est bijectif.

\uline{Preuve}. Soit $N$ un entier suffisamment grand pour que
\[
\bigcup_{n\underline{>}0} \mathrm{Ker}\, V^{n} = \mathrm{Ker}\, V^{N} \quad \text{et} \quad \bigcap_{n\underline{>}0} \mathrm{Im}\, V^{n} = \mathrm{Im}\, V^{N} .
\]
On a évidemment $V|\mathrm{Ker}\, V^{N}$ nilpotent et $V|\mathrm{Im}\, V^{N}$ bijectif et de plus
$M = \mathrm{Ker}\, V^{N} \oplus \mathrm{Im}\, V^{N}$ . Il est clair que cette décomposition est unique et
fonctorielle.

5. \uline{Corollaires et compléments}

\uline{Corollaire 5.1}

Le foncteur $D^*$ commute aux changements de base définis par
toute extension parfaite $K$ de $k$ : Pour un objet $G$ de $p$-Gr-fin$(k)$ ,
on a un isomorphisme fonctoriel
\[
D^*(G_{K}) \overset{\sim}{\leftarrow} W(K) \otimes_{W(k)} D^*(G) .
\]

\mspage{\hand{47}}

\clearpage
% Notes of the Montreal seminar, 1970 (cote 129, typed part): archive cote 129, site page 48 (PDF page 49).
% typescript. Draft compared: none.
% AI-transcribed by Claude Opus 5.5 (claude-opus-5-5, adaptive thinking, effort high), Message Batches API, batch mtl129_full (msgbatch_01AGnDKHwy6ATCM7ehUa6263), 2026-11-13; not yet checked by eye.
\scanpage{129/48}
\mspage{\hand{II 20}}

\uline{Preuve} (d'après Oda et Oort).

On a une flèche évidente. Il suffit de considérer le cas où
$G$ est unipotent, et en considérant alors la suite de composition formée
des $\mathrm{Im}\,\mathbb{V}^n_G \subset G$ , on se ramène au cas où $\mathbb{V}_G = 0$ .

Pour un tel groupe $G$ , on a
\[ D^*(G) = \mathrm{Hom}_{k\text{-}Gr}(G,\underset{\to}{\mathbb{W}}) = \mathrm{Hom}_{k\text{-}Gr}(G,\mathbb{G}_a) . \]

Par ailleurs, on a un isomorphisme
\[ \mathrm{hom}_{k\text{-}Gr}(G,\mathbb{G}_a) \cong \mathrm{Ker}\{H^0(G,\underline{O}_G) \xrightarrow{\ \Delta^*-pr_1^*-pr_2^*\ } H^0(G,\underline{O}_G) \otimes_k H^0(G,\underline{O}_G)\} \]
où $\Delta : G \times G \to G$ désigne la multiplication de $G$ . Par le changement de
base $k \to K$ , on a un isomorphisme analogue, d'où l'on déduit le corol-%
laire.

\uline{Corollaire 5.2}

Pour tout p-groupe fini sur $k$ , on a
\[ \mathrm{rang}(G) = p^{\mathrm{long}_W(D^*(G))} . \]

\uline{Preuve}. Il suffit de considérer le cas où $G$ est unipotent, et par le
Corollaire 5.1 on peut supposer $k$ algébriquement clos. Alors $G$ admet
une suite de composition ne faisant intervenir que $\mathbb{Z}/p\mathbb{Z}$ et $\alpha_p$ (2.4).
Or on établit aisément que
\[ D^*(\alpha_p) \sim k = W/pW \quad \text{avec} \quad F = V = 0 \]
\[ D^*(\mathbb{Z}/p\mathbb{Z}) \sim k = W/pW \quad \text{avec} \quad V = 0 \ , \ F(\lambda) = \lambda p \ , \]
et la formule est trivialement vérifiée.

\margin{48}

\clearpage
% Notes of the Montreal seminar, 1970 (cote 129, typed part): archive cote 129, site page 49 (PDF page 50).
% typescript. Draft compared: none.
% AI-transcribed by Claude Opus 5.5 (claude-opus-5-5, adaptive thinking, effort high), Message Batches API, batch mtl129_full (msgbatch_01AGnDKHwy6ATCM7ehUa6263), 2026-11-13; not yet checked by eye.
\scanpage{129/49}
\mspage{\hand{II 21}}

5.3. \uline{Remarques}

5.3.1. Les groupes formels et les groupes de Barsotti-Tate sur $k$
peuvent être considérés comme limites inductives de p-groupes finis sur
$k$ . La théorie de Dieudonné permet donc de décrire ces groupes en termes
de limites projectives de $\mathcal{D}$-modules de longueur finie sur $W$ . Par exem-%
ple, on montrera (III, 5.6) que l'on a une équivalence de catégories en-%
tre la catégorie des groupes de Barsotti-Tate sur $k$ et la catégorie
des $\mathcal{D}$-modules qui sont libres et de type fini sur $W$ .

5.3.2. Pour un p-groupe $G$ fini sur $k$ , il existe un isomorphisme
canonique de $\mathcal{D}$-modules, fonctoriel en $G$ :
\[ D^*(G^*) \overset{\sim}{\to} \mathrm{Hom}_W(D^*(G),\ K/W) \ , \]
où $K/W$ désigne le module dualisant de $W$ (c.f. Gabriel et Demazure).
La structure de $\mathcal{D}$-module sur le membre de droite est définie comme suit :

Pour tout $\alpha \in \mathrm{Hom}_W(D^*(G),\ K/W)$ et tout $x \in D^*(G)$ , on pose
\[ F{\cdot}\alpha(x) = [\alpha(V(x))]^{\sigma} \]
\[ V{\cdot}\alpha(x) = [\alpha(F(x))]^{\sigma-1} \ . \]

6. \uline{Annexe = Construction du foncteur quasi inverse de} $D^*$

Il suffit de construire le foncteur quasi-inverse du foncteur
\begin{diagramblock}
\[ D^* : \quad \left(\begin{array}{c} \text{groupes finis} \\ \text{unipotents} \end{array}\right)^{0} \longrightarrow \left(\begin{array}{l} \text{sous catégorie de } \underline{\mathrm{Mod}}_{\mathcal{D};\,W\ \hand{\ell.f.}} \text{ for-} \\ \text{mée des objets où } V \text{ est nilpotent} \end{array}\right) \]
\end{diagramblock}
défini en 4.3. On montrera que le foncteur que l'on va construire est un

\mspage{\hand{49}}

\clearpage
% Notes of the Montreal seminar, 1970 (cote 129, typed part): archive cote 129, site page 50 (PDF page 51).
% typescript. Draft compared: none.
% AI-transcribed by Claude Opus 5.5 (claude-opus-5-5, adaptive thinking, effort high), Message Batches API, batch mtl129_full (msgbatch_01AGnDKHwy6ATCM7ehUa6263), 2026-11-13; not yet checked by eye.
\scanpage{129/50}
\mspage{\hand{II.22}}

adjoint à gauche de $D^*$ , il résultera du fait que $D^*$ est une équivalen-%
ce de catégories que c'est un foncteur quasi-inverse de $D^*$ .

6.1. Soit $M$ un $\hand{\mathcal{D}}$-module de longueur finie sur $W$ et où $V$ est
nilpotent sur $M$ . On lui associe le foncteur en groupes $\mathbb{E}(M)$ sur
(sch/k) défini par
\[ \mathbb{E}(M)(S) = \mathrm{Hom}_{\hand{\mathcal{D}}}(M, \underset{\hand{\to}}{\mathbb{W}}(S)) \]
pour tout k-schéma $S$ .

Il faut démontrer que $\mathbb{E}(M)$ est représenté par un p-groupe
fini unipotent et que le foncteur $\mathbb{E}$ ainsi défini est adjoint à gauche
de $D^*$ .

6.2. $\mathbb{E}(M)$ \uline{est représentable par un schéma affine}

On choisit un entier $N$ suffisamment grand pour que $V^{N+1}$
soit nul sur $M$ et on considère les polynômes $P_N$ et $S_N$ qui définis-%
sent respectivement la dernière coordonnée de la multiplication et de
l'addition dans $\mathbb{W}_{N+1}$ (ce sont des polynômes sur $\mathbb{Z}$ à $2N+2$ variables).

On définit alors la k-algèbre $A$ comme quotient de
$k[T_x, x \in M]$ par les relations suivantes :
\[ \text{i)}\quad T_{F(x)} = T_x^p \]
\[ \text{ii)}\quad T_{x+y} = S_N(T_{V^N(x)}, \ldots, T_x \,;\, T_{V^N(y)}, \ldots, T_y) \]
\margin{$\times$}
\[ \text{iii)}\quad T_{\hand{\lambda}x} = P_N(\hand{\lambda}_1^{p^{1-N}}, \hand{\lambda}_2^{p^{1-N}}, \ldots, \hand{\lambda}_{N+1}^{p^{1-N}} \,;\, T_{V^N(\hand{x})}, \ldots, T_{\hand{x}}) \ . \]

Pour prouver que $\mathrm{Spec}(A)$ représente le foncteur $\mathbb{E}(M)$ ,
on remarque que pour tout $x \in M$ et tout $u \in \mathrm{Hom}_{\hand{\mathcal{D}}}(M, \underset{\hand{\to}}{\mathbb{W}}(S))$ , on a

\mspage{\hand{50}}

\clearpage
% Notes of the Montreal seminar, 1970 (cote 129, typed part): archive cote 129, site page 51 (PDF page 52).
% typescript. Draft compared: none.
% AI-transcribed by Claude Opus 5.5 (claude-opus-5-5, adaptive thinking, effort high), Message Batches API, batch mtl129_full (msgbatch_01AGnDKHwy6ATCM7ehUa6263), 2026-11-13; not yet checked by eye.
\scanpage{129/51}
\mspage{\hand{II 23}}

$V^{N+1}(u(x)) = 0$ car $V^{N+1}(x) = 0$ ; l'élément $u(x)$ provient par la
limite inductive d'un élément $(y_1,y_2,\ldots,y_r) \in \mathbb{W}_r(S)$ et on a
$V^{N+1}(y_1,\ldots,y_r) = 0$ car $\mathbb{W}_r$ se plonge dans $\underset{\rightarrow}{\mathbb{W}}$ ; on a donc
$y_1 = y_2 = \ldots y_{r-N-1} = 0$ , ce qui implique que $(y_1,y_2,\ldots,y_r)$ pro-%
vient d'un élément unique de $\mathbb{W}_{N+1}(S)$ par le morphisme de transition
$T^{r-N-1}$ .

On en déduit que tout homomorphisme $u \in \mathrm{Hom}_{\mathcal{D}}(M, \mathbb{W}(S))$ se
factorise de manière unique par $W_{N+1}(S) \hookrightarrow \underset{\rightarrow}{\mathbb{W}}(S)$ .

On est donc ramené au problème de construire un isomor-%
phisme fonctoriel en $S$
\[ \phi_S : \mathrm{Hom}_{\mathcal{D}}(M, \mathbb{W}_{N+1}(S)) \rightarrow \mathrm{Hom}_k(A, \Gamma(\underline{\mathcal{O}}_S)) . \]

Pour tout $u \in \mathrm{Hom}_{\mathcal{D}}(M, \mathbb{W}_{N+1}(S))$ , $\phi_S(u) : A \rightarrow \Gamma(\mathcal{O}_S)$ est
le morphisme qui associe à $T_x$ la dernière coordonnée de
\margin{$\times$} $u(x) \in \mathbb{W}_{N+1}(S) = \Gamma(\mathcal{O}_S)^{N+1}$ . Il est clair que $\hand{\phi}_S(u)$ est bien défini
sur $A$ et est un morphisme d'algèbres (à cause des relations vérifiées
dans $A$). On définit une application inverse en associant à $\overline{u} \in \mathrm{Hom}_k(A, \Gamma(\mathcal{O}_S))$
l'homomorphisme de $\mathcal{D}$-modules $u : M \rightarrow W_{N+1}(S)$ qui à $x$ fait corres-%
pondre $u(x) = (\overline{u}(T_{V^N(x)})\ldots, \overline{u}(T_x))$ .

6.3. \uline{Le schéma} $\mathbb{E}(M)$ \uline{est un p-groupe fini unipotent}

Par définition même de $\mathbb{E}(M)$ (6.1), $\mathbb{E}(M)$ est un schéma
en groupes, dépendant additivement de $M$ . Comme $M$ est annulé par une
puissance de $p$ (car $V$ est nilpotent), il en est de même de $\mathbb{E}(M)$ .
Montrons que $\mathbb{E}(M)$ est fini.

\mspage{\hand{51}}

\clearpage
% Notes of the Montreal seminar, 1970 (cote 129, typed part): archive cote 129, site page 52 (PDF page 53).
% typescript. Draft compared: none.
% AI-transcribed by Claude Opus 5.5 (claude-opus-5-5, adaptive thinking, effort high), Message Batches API, batch mtl129_full (msgbatch_01AGnDKHwy6ATCM7ehUa6263), 2026-11-13; not yet checked by eye.
\scanpage{129/52}
\mspage{\hand{II 24}}

Soit $\{x_1,\dots,x_r\}$ un système de générateurs du W-module $M$ ;
les relations vérifiées dans l'anneau $A$ entraînent que l'on peut trou-%
ver un système de monômes en les variables $T_{x_1},\dots,T_{x_r}$ , $T_{Vx_1},\dots,T_{Vx_r},\dots,$
$T_{V^Nx_1},\dots,T_{V^Nx_r}$ et de degrés bornés, tels que cette famille de monômes
engendre $A$ en tant que k-espace vectoriel.

D'après 2.5.2 il suffit de montrer que $\mathbb{V}_{\mathbb{E}(M)} = \mathbb{V}$ est nilpo-%
tent pour prouver que $\mathbb{E}(M)$ est unipotent. Comme $\mathbb{E}(M)$ est un sous-%
groupe d'un groupe lisse (l'espace affine de dimension $r(N+1)$) , $v$ est
déterminé de manière unique (I.5.3) par la relation $\mathbb{V}\circ\mathbb{F}_{\mathbb{E}(M)|k} = p$ ;
nous allons exhiber un morphisme $\mathbb{V}$ satisfaisant à cette relation. On
montre que l'on a le diagramme commutatif suivant

\begin{diagramblock}
\begin{tikzcd}[column sep=large]
\mathrm{Hom}_{\mathcal{D}}(M, \mathbb{W}_{n+1}(S)) \arrow[r, "\alpha \to \alpha\circ F"] \arrow[d, "\phi_S", "\wr"'] & \mathrm{Hom}_{\mathcal{D}}(M, \mathbb{W}_{N+1}(S,\sigma)) \arrow[d, "\phi_{S,\sigma}", "\wr"'] \\
\mathrm{Hom}_{k\text{-alg}}(A, \Gamma(\underline{\mathcal{O}}_S)) \arrow[r, "\mathbb{F}_{\mathbb{E}(M)|k}(S)"] & \mathrm{Hom}_{k\text{-alg}}(A \otimes_k (k,\sigma), \Gamma(\underline{\mathcal{O}}_S))
\end{tikzcd}
\end{diagramblock}

où $(S,\sigma)$ désigne le k-schéma $S \to \mathrm{Spec}(k) \overset{\sigma}{\to} \mathrm{Spec}(k)$ .

On construit alors un morphisme fonctoriel en $S$
\[ \mathbb{V} : \mathrm{Hom}_{\mathcal{D}}(M,\mathbb{W}_{N+1}(S,\sigma)) \to \mathrm{Hom}_{\mathcal{D}}(M,\mathbb{W}_{N+1}(S)) \]
\[ u \to u \circ V \ . \]

D'après ce qui a été dit, ce morphisme s'identifie à la Vers-%
chiebung. Comme $V$ est nilpotent sur $M$ , $\mathbb{V}$ est nilpotent et $\mathbb{E}(M)$
est unipotent.

\margin{\hand{52}}

\clearpage
% Notes of the Montreal seminar, 1970 (cote 129, typed part): archive cote 129, site page 53 (PDF page 54).
% typescript. Draft compared: none.
% AI-transcribed by Claude Opus 5.5 (claude-opus-5-5, adaptive thinking, effort high), Message Batches API, batch mtl129_full (msgbatch_01AGnDKHwy6ATCM7ehUa6263), 2026-11-13; not yet checked by eye.
\scanpage{129/53}
\mspage{\hand{II.25}}

6.4. \uline{Le foncteur} $\mathbb{E}$ \uline{est adjoint à gauche de} D*

On veut construire un isomorphisme fonctoriel en G et M
\[ \mathrm{Hom}_{k\text{-}Gr}(G,\mathbb{E}(M)) \overset{\sim}{\to} \mathrm{Hom}_{\hand{$\mathcal{D}$}}(M,\ D^*(G)) \ . \]
Si l'on pose
\margin{XX}
\[ \mathrm{Prim}(\deleted{(}G) = \mathrm{Ker}\ \{\underset{\rightarrow}{\mathbb{W}}(G) \xrightarrow{\ \Delta^*-pr_1^*-pr_2^*\ } \underset{\rightarrow}{\mathbb{W}}(G \times G)\}\deleted{)}\ , \]
on vérifie que l'on a :
\[ D^*(G) = \mathrm{Hom}_{k\text{-}Gr}(G,\underset{\rightarrow}{\mathbb{W}}) \overset{\sim}{\to} \mathrm{Prim}(G)\ . \]
Par ailleurs, d'après la définition même de $\mathbb{E}(M)$ , on a
\[ \mathrm{Hom}_k(G,\ \mathbb{E}(M)) = \mathbb{E}(M)(G) = \mathrm{Hom}_{\hand{$\mathcal{D}$}}(M,\ \underset{\rightarrow}{\mathbb{W}}(G))\ , \]
ce qui définira l'isomorphisme désiré si l'on montre qu'un élément
$u \in \mathrm{Hom}_{\hand{$\mathcal{D}$}}(M,\ \underset{\rightarrow}{\mathbb{W}}(G))$ définit un morphisme de groupes $G \to \mathbb{E}(M)$ si et seu-%
lement si u se factorise par Prim(G) .

En effet $\mathrm{Hom}_{k\text{-}Gr}(G,\ \mathbb{E}(M))$ s'identifie à la partie de
$\mathrm{Hom}_{k\text{-}alg}(A,\ \Gamma(\hand{$\underline{\mathcal{O}}$}_G))$ formée des morphismes $\alpha$ tels que le diagramme sui-%
vant soit commutatif :
\begin{diagramblock}
\begin{tikzcd}
A \arrow[r, "\alpha"] \arrow[d, "\Delta"'] & \Gamma(\hand{$\underline{\mathcal{O}}$}_G) \arrow[d, "\Delta"] \\
A \otimes A \arrow[r, "\alpha \otimes \alpha"] & \Gamma(\hand{$\underline{\mathcal{O}}$}_G) \otimes \Gamma(\hand{$\underline{\mathcal{O}}$}_G)\ .
\end{tikzcd}
\end{diagramblock}

Mais le morphisme $\Delta : A \to A \otimes A$ est donné par la formule
\[ \Delta(T_x) = S_N(T_{V^N(x)} \otimes 1, \ldots, T_x \otimes 1\ ;\ 1 \otimes T_{V^N(x)}, \ldots,\ 1 \otimes T_x) \]

\margin{53}

\clearpage
% Notes of the Montreal seminar, 1970 (cote 129, typed part): archive cote 129, site page 54 (PDF page 55).
% typescript. Draft compared: none.
% AI-transcribed by Claude Opus 5.5 (claude-opus-5-5, adaptive thinking, effort high), Message Batches API, batch mtl129_full (msgbatch_01AGnDKHwy6ATCM7ehUa6263), 2026-11-13; not yet checked by eye.
\scanpage{129/54}
\mspage{\hand{II.26}}

pour tout \ $x \in M$ , car pour \ $u$ , $u' \in \mathrm{Hom}_{\mathcal{D}}(M, \mathbb{W}_{N+1}(S))$ , on a
$\phi_S(u + u')(T_x) = S_N(u(x), u'(x))$ \ et, toujours d'après la définition
de \ $\phi_S$ , on a, en posant \ $\overline{u} = \phi_S(u)$ , $\overline{u}' = \phi_S(u')$
\[ S_N(u(x),u'(x)) = S_N[\overline{u}(T_{V^N(x)}),\dots,\overline{u}(T_x) \;;\; \overline{u}'(T_{V^N(x)}),\dots,\overline{u}'(T_x)] \ . \]

De sorte que \ $\overline{u}$ \ est un morphisme de groupe si et seule-%
ment si pour tout \ $x \in M$ , on a
\[ \Delta(\overline{u}(T_x)) = S_N[\overline{u}(T_{V^N(x)}) \otimes 1,\dots,\overline{u}(T_x) \otimes 1;\ 1 \otimes \overline{u}(T_{V^N(x)}),\dots,1 \otimes \overline{u}(T_x)] \ . \]

Ce dernier élément n'est autre que la N+1$^{\text{ième}}$ coordonnée
\margin{$\times$ \ $\otimes$} dans \ $\mathbb{W}_{N+1}(\Gamma(\underline{\mathcal{O}}_G)$ \hand{(}$x$\hand{)} $\Gamma(\underline{\mathcal{O}}_G))$ \ de l'élément \ $u(x) \otimes 1 + 1 \otimes u(x)$ .

Le même raisonnement appliqué à \ $x$ , $V(x),\dots,V^N(x)$ \ montre
alors que \ $\overline{u}$ \ définit un morphisme de groupes \ $G \to \mathbb{E}(M)$ \ si et seulement
si \ $u(x)$ \ appartient au noyau de \ $\mathbb{W}_{N+1}(\Gamma(\underline{\mathcal{O}}_G)) \xrightarrow{\ \Delta^*\text{-}\mathrm{pr}_1^*\text{-}\mathrm{pr}_2^*\ } \mathbb{W}_{N+1}(\Gamma(\underline{\mathcal{O}}_G) \otimes \Gamma(\underline{\mathcal{O}}_G))$ ,
qui est lui-même contenu dans \ Prim(G) .

\margin{54}

\clearpage
% Notes of the Montreal seminar, 1970 (cote 129, typed part): archive cote 129, site page 55 (PDF page 56).
% typescript. Draft compared: none.
% AI-transcribed by Claude Opus 5.5 (claude-opus-5-5, adaptive thinking, effort high), Message Batches API, batch mtl129_full (msgbatch_01AGnDKHwy6ATCM7ehUa6263), 2026-11-13; not yet checked by eye.
\scanpage{129/55}
CHAPITRE III

GROUPES DE BARSOTTI-TATE

\mspage{\hand{55}}

\clearpage
% Notes of the Montreal seminar, 1970 (cote 129, typed part): archive cote 129, site page 56 (PDF page 57).
% typescript. Draft compared: none.
% AI-transcribed by Claude Opus 5.5 (claude-opus-5-5, adaptive thinking, effort high), Message Batches API, batch mtl129_full (msgbatch_01AGnDKHwy6ATCM7ehUa6263), 2026-11-13; not yet checked by eye.
\scanpage{129/56}
\mspage{\hand{III-1}}

\uline{GROUPES DE BARSOTTI-TATE}

1. \uline{Notations et définitions préléminaires}

Soit $S$ un schéma. On munit la catégorie Sch/S des schémas
au dessus de $S$ d'une topologie \hand{$\mathcal{T}$} que l'on suppose plus fine que la
topologie f.p.p.f. (fidèlement plate et de présentation finie) et moins
fine que la topologie f.p.q.c. (fidèlement plate et quasi-compacte)
(SGA 3 IV, 6.3). L'introduction de cette topologie est technique et n'est
pas essentielle dans les énoncés finaux (voir 5.1).

Pour abréger, nous dirons ici que $G$ est un groupe si $G$ est
un faisceau en groupes commutatifs sur Sch/S pour la topologie \hand{$\mathcal{T}$}.

Pour tout entier $n \geq 1$ et pour tout groupe $G$, on pose:
\[ \Lambda_n = \mathbb{Z}/p^n\mathbb{Z} \]
\[ G(n) = \mathrm{Ker}\,\{G \xrightarrow{p^n \mathrm{Id}} G\} \ , \]
et
\[ G(\infty) = \varinjlim_{n} G(n) \ ; \]
$G(n)$ est un $\Lambda_n$-module (i.e. un \hand{$\mathcal{T}$}-faisceau en $\Lambda_n$-modules);
$G(\infty)$ est le plus grand sous-groupe de $p$-torsion de $G$. On a évidemment:
\[ G(n)(n') = G(n') \qquad \text{pour} \quad n' \leq n \leq \infty \]

Si l'on veut indiquer qu'un groupe est annulé par $p^n$, c'est-à-%
dire est un $\Lambda_n$-module, on pourra donc le noter $G(n)$ au lieu de $G$ sans
entraîner de confusions.

\mspage{\hand{56}}

\clearpage
% Notes of the Montreal seminar, 1970 (cote 129, typed part): archive cote 129, site page 57 (PDF page 58).
% typescript. Draft compared: none.
% AI-transcribed by Claude Opus 5.5 (claude-opus-5-5, adaptive thinking, effort high), Message Batches API, batch mtl129_full (msgbatch_01AGnDKHwy6ATCM7ehUa6263), 2026-11-13; not yet checked by eye.
\scanpage{129/57}
\mspage{\hand{III-2}}

2. \uline{Platitude et critère de représentabilité}

2.1. Soit $G$ un $\Lambda_n$-module. On dit que $G$ est \uline{plat sur} $\Lambda_n$ si
le produit tensoriel par $G$ au dessus du faisceau constant $\Lambda_n$ préserve
les monomorphismes; on a alors le critère suivant:

2.2. \uline{Proposition}: Soit $G(n)$ un groupe annulé par $p^n$. On a les
équivalences:

(a) $\forall i \quad 1 \leq i \leq n, \quad G(i)$ est un $\Lambda_i$-module plat

(b) $G(n)$ est un $\Lambda_n$-module plat

(c) $\forall i, \quad 1 \leq i \leq n-1, \quad p^{n-i}\, G(n) = G(i)$

(i.e. l'application $G(n) \xrightarrow{p^{n-i}\mathrm{Id}} G(i)$ est un épimorphisme)

(d) $\exists i, \quad 1 \leq i \leq n-1,$ tel que $p^{n-i}G(n) = G(i)$

(e) $p\, G(n) = G(n-1)$

(f) Le morphisme canonique
\[ \gamma\colon \Lambda_1[T]/T^n \underset{\Lambda_1}{\otimes} G(n)/p\, G(n) \to gr_p G(n) = \bigoplus_{i=0}^{i=n-1} p^i G(n)/p^{i+1}G(n) \]
(où $\gamma^i\colon G(n)/p\, G(n) \to p^i G(n)/p^{i+1}G(n)$ est induit par la
multiplication par $p^i$) est un isomorphisme.

(g) Le morphisme
\[ \beta\colon gr_p G(n) \to \Lambda_1[T]/T^n \underset{\Lambda_1}{\otimes} G(1) \]
(où $\beta^i\colon p^i G(n)/p^{i+1}G(n) \to G(1)$ est induit par la multipli-%
cation par $p^{n-i-1}$) est un isomorphisme.

\uline{Démonstration}:

(b) $<=>$ (c) Pour tout $i$, $1 \leq i \leq n-1$, on a des triangles commutatifs

\mspage{\hand{57}}

\clearpage
% Notes of the Montreal seminar, 1970 (cote 129, typed part): archive cote 129, site page 58 (PDF page 59).
% typescript. Draft compared: none.
% AI-transcribed by Claude Opus 5.5 (claude-opus-5-5, adaptive thinking, effort high), Message Batches API, batch mtl129_full (msgbatch_01AGnDKHwy6ATCM7ehUa6263), 2026-11-13; not yet checked by eye.
\scanpage{129/58}
\mspage{III-3}

canoniques de $\Lambda_n$-modules constants qui donnent, par produit tensoriel
avec $G(n)$ les diagrammes suivants:

\begin{diagramblock}
\begin{tikzcd}
p^i\Lambda_n \arrow[r, hook, "\varepsilon"] & \Lambda_n \\
\Lambda_n/p^{n-i}\Lambda_n \arrow[u, "p^i", "\wr"'] \arrow[ur, "(p^i)"'] &
\end{tikzcd}
\qquad
\begin{tikzcd}
p^i\Lambda_n\otimes G(n) \arrow[rr, "\varepsilon\otimes 1"] & & G(n) \\
\Lambda_n/p^i\Lambda_n\otimes G(n) = G(n)/p^iG(n) \arrow[u, equal, "\wr"'] \arrow[urr, "(p^i)"'] & &
\end{tikzcd}
\end{diagramblock}

les seuls idéaux de $\Lambda_n$ étant les $p^i\Lambda_n$, $G(n)$ est plat si et seulement
si $\varepsilon\otimes 1$ est injectif pour tout $i$, ou encore si et seulement si $(p^i)$
est injectif, ce qui équivaut à (c).

(c) $\Rightarrow$ (a) Les relations analogues à (c) sont alors vérifiées par
tous les $G(i)$ qui sont donc plats sur $\Lambda_i$, d'après ce qui précède.

(d) $\Rightarrow$ (e) Montrons que si $G(i) = p^{n-i}G(n)$, pour un certain $i$,
$1 \leq i \leq n-1$ alors on a $G(i+1) = p^{n-i-1}G(n)$. Cela se voit directement sur
le diagramme suivant où les trois lignes sont des suites exactes.

\begin{diagramblock}
\begin{tikzcd}
0 \arrow[r] & G(n-1) \arrow[r, hook] \arrow[d, "p^{n-i-1}"] & G(n) \arrow[r, "p^{n-1}"] \arrow[d, "p^{n-i-1}"] & p^{n-1}G = p^{i-1}G(i) \arrow[r] \arrow[d, equal] & 0 \\
0 \arrow[r] & G(i) \arrow[r] \arrow[d] & G(i+1) \arrow[r, "p^i"] \arrow[d] & p^iG(i+1) = p^{i-1}G(i) \arrow[r] \arrow[d] & 0 \\
 & 0 \arrow[r, twoheadrightarrow] \arrow[d] & G(i+1)/p^{n-i-1}G(n) \arrow[r] \arrow[d] & 0 \arrow[r] \arrow[d] & 0 \\
 & 0 & 0 & 0 &
\end{tikzcd}
\end{diagramblock}

(e) $\Leftrightarrow$ (f) Soit $\lambda^i : p^iG(n)/p^{i+1}G(n) \rightarrow p^{i+1}G(n)/p^{i+2}G(n)$, $0 \leq i \leq n-2$

\mspage{58}

\clearpage
% Notes of the Montreal seminar, 1970 (cote 129, typed part): archive cote 129, site page 59 (PDF page 60).
% typescript. Draft compared: none.
% AI-transcribed by Claude Opus 5.5 (claude-opus-5-5, adaptive thinking, effort high), Message Batches API, batch mtl129_full (msgbatch_01AGnDKHwy6ATCM7ehUa6263), 2026-11-13; not yet checked by eye.
\scanpage{129/59}
\mspage{III-4}

l'épimorphisme induit par la multiplication par $p$. On a
\[ \gamma^{i} = \lambda^{i-1} \circ \lambda^{i-2} \circ \ldots \circ \lambda^{o}, \quad 1 \leq i \leq n-1 \]
on en déduit que l'épimorphisme $\gamma$ est un isomorphisme si et seulement
si $\gamma^{n-1}$ est un monomorphisme, i.e. si on a (e).

(e) $\Rightarrow$ (g) $\Rightarrow$ (c) Avec les notations précédentes, on a:
\[ \beta^{i} = \deleted{J}\inserted{j} \circ \lambda^{n-2} \circ \lambda^{n-3} \circ \ldots \circ \lambda^{i} \] \margin{x \quad j}
(où $j$ désigne l'injection canonique $p^{n-1}G(n) \rightarrow G(1)$). D'après ce qui
précède, (e) entraîne que tous les $\lambda^{i}$ sont des isomorphismes, i.e.
$\beta$ est un monomorphisme. Nous allons montrer par récurrence que si $\beta$ est
un monomorphisme, on a (c), d'où l'on déduit immédiatement que $j$ et donc
$\beta$, sont des isomorphismes.

Si $\beta^{o} = j\circ\gamma^{n-1}$ est un monomorphisme, on a alors $pG(n) = G(n-1)$.
Comme $\beta^{i}$ est un monomorphisme, on a:
\[ [p^{i}G(n)](n-i-1) = p^{i+1}G(n) \ ; \]
mais $p^{i}G(n) = G(n-i)$ par l'hypothèse de récurrence, et donc
$G(n-i-1) = p^{i+1}G(n)$.

2.3. Supposons maintenant que $G(n)$ soit un $\Lambda_{n}$-module plat et qu'il
soit représentable par un schéma fini localement libre sur $S$. Alors $G(i)$,
noyau de la multiplication par $p^{i}$, est représentable, fini et localement
de présentation finie. Mais, $G(n)$ étant supposé $\Lambda_{n}$-plat, on a un épi-%
morphisme $G(n) \overset{p^{n-1}}{\twoheadrightarrow} G(i)$. Fibre à fibre, cette application est surjective

\mspage{59}

\clearpage
% Notes of the Montreal seminar, 1970 (cote 129, typed part): archive cote 129, site page 60 (PDF page 61).
% typescript. Draft compared: none.
% AI-transcribed by Claude Opus 5.5 (claude-opus-5-5, adaptive thinking, effort high), Message Batches API, batch mtl129_full (msgbatch_01AGnDKHwy6ATCM7ehUa6263), 2026-11-13; not yet checked by eye.
\scanpage{129/60}
\mspage{\hand{III-5}}

et donc (fidèlement) plate.

Par le critère de platitude fibre à fibre (EGA, IV, 11.3.1)
on en déduit que $G(i)$ est plat sur $S$. Alors $G(i)$ est plat, fini,
localement de présentation finie sur $S$, et est donc localement libre.

Inversement, si $G(1)$ est représentable par un schéma fini
localement libre sur $S$ et si $G(n)$ est $\Lambda_n$-plat, alors $G(n)$ admet
une suite de décomposition finie $p^iG(n)/p^{i+1}G(n) \approx G(1)$ et $G(n)$ est
représentable par un schéma fini localement libre (II, 1.6). On a donc:

2.4. \uline{Proposition}: Soit $G(n)$ un groupe annulé par $p^n$ et plat sur
$\Lambda_n$. Alors $G(n)$ est représentable par un schéma fini localement libre
sur $S$ si et seulement si l'un des $G(i)$, $1 \leq i \leq n$, l'est.

3. \uline{Groupes de Barsotti-Tate tronqués d'échelon $n$}

3.1. Nous allons voir que la définition d'un groupe de Barsotti-Tate
tronqué d'échelon $n$ fait intervenir la platitude sur $\Lambda_n$. Pour $n = 1$
cette condition est vide et il est nécessaire d'introduire une condition
supplémentaire que l'on justifiera a posteriori (3.3).

Soit $G$ un schéma en groupes (commutatifs) plat sur $S$. On
considère le $\mathbb{F}_p$-schéma fermé $S_o \subset S$ défini par l'idéal $p\underline{\mathcal{O}}_S$ et soit
$G_o$ le schéma sur $S_o$ obtenu par changement de base. On peut alors cons-%
truire les morphismes $\mathfrak{f}_{G_o/S_o}$ et $\mathfrak{v}_{G_o}$. Si $G_o$ est un $\Lambda_1$-module, on a
$\mathfrak{f}_{G_o/S_o} \circ \mathfrak{v}_{G_o} = 0$.

\mspage{\hand{60}}

\clearpage
% Notes of the Montreal seminar, 1970 (cote 129, typed part): archive cote 129, site page 61 (PDF page 62).
% typescript. Draft compared: none.
% AI-transcribed by Claude Opus 5.5 (claude-opus-5-5, adaptive thinking, effort high), Message Batches API, batch mtl129_full (msgbatch_01AGnDKHwy6ATCM7ehUa6263), 2026-11-13; not yet checked by eye.
\scanpage{129/61}
\mspage{III-6}

3.2. \uline{Définition}. Soit $n \geq 1$ un entier. On appelle \uline{p-groupe de}
\uline{Barsotti-Tate tronqué d'échelon $n$ sur $S$} (ou groupe Barsotti-Tate
tronqué d'échelon $n$) un schéma en groupes (commutatifs) $G$ fini et
localement libre sur $S$, annulé par $p^{n}$ et plat sur $\Lambda_{n}$. Si $n = 1$,
on suppose de plus que l'on a (avec les notations de 3.1):
\[
(*) \qquad \operatorname{Ker} \mathbb{F}_{G_{o}/S_{o}} = \operatorname{Im} \mathbb{V}_{G_{o}}
\]
(c'est une condition qui se vérifie fibre à fibre).

3.3. Si $G(n)$ est un groupe de BT tronqué d'échelon $n$, alors
$G(i)$, $1 \leq i \leq n$, est un groupe de BT tronqué d'échelon $i$.

Pour $i > 1$, cela résulte de 2.4; pour $i = 1$, nous allons
montrer que la condition (*) est aussi vérifiée.

Au dessus de $S_{o} \subset S$, nous avons déjà remarqué que l'image de
$\mathbb{V}_{G(1)_{o}} : G(1)_{o}^{(p)} \rightarrow G(1)_{o}$ était contenue dans le noyau $K$ de
$\mathbb{F}_{G(1)_{o}/S_{o}} : G(1)_{o} \rightarrow G(1)_{o}^{(p)}$ , de sorte que l'on a un diagramme commutatif
\begin{diagramblock}
\begin{tikzcd}
G(1)_{o}^{(p)} \arrow[r, "\mathbb{V}"] \arrow[dd] & K \arrow[d, hook] \\
 & G(1)_{o} \arrow[d, hook] \\
G(2)_{o}^{(p)} \arrow[r, "\mathbb{V}_{G(2)_{o}}"] & G(2)_{o}
\end{tikzcd}
\end{diagramblock}
dont on vérifie immédiatement qu'il est cartésien car $G(2)_{o}^{(p)}(1) \simeq G(1)_{o}^{(p)}$.
Par ailleurs l'image de $\mathbb{V}_{G(2)_{o}}$ contient $G(1)_{o}$ , car
$\mathbb{V}_{G(2)_{o}} \circ \mathbb{F}_{G(2)_{o}} = p$ a pour image $G(1)_{o}$ , et l'on en déduit que
$\mathbb{V} : G(1)_{o}^{(p)} \rightarrow K$ est surjectif.

\mspage{61}

\clearpage
% Notes of the Montreal seminar, 1970 (cote 129, typed part): archive cote 129, site page 62 (PDF page 63).
% typescript. Draft compared: none.
% AI-transcribed by Claude Opus 5.5 (claude-opus-5-5, adaptive thinking, effort high), Message Batches API, batch mtl129_full (msgbatch_01AGnDKHwy6ATCM7ehUa6263), 2026-11-13; not yet checked by eye.
\scanpage{129/62}
\mspage{\hand{III-7}}

4. \uline{Groupes de Barsotti-Tate}

La terminologie de groupe de Barsotti-Tate est préférée à
celle de ``groupe p-divisible'' employée par Tate (J.T. Tate, p-divisi-%
ble groups, in Proceedings of a Conference on Local Fields, Driebergen
1966 (Springer)), la notion de groupe p-divisible étant beaucoup plus
générale:

4.1. \uline{Définition}. On dit qu'un groupe $G$ ($\mathcal{T}$-faisceau en groupes)
est p-divisible si $p.\mathrm{Id}\colon G \to G$ est un épimorphisme.

Si $G$ est p-divisible, $G(n)$ est un $\Lambda_n$-module plat d'après
2.2 (e), et le groupe de p-torsion $G(\infty)$ est lui aussi p-divisible.
Plus généralement, si on se donne un système inductif
\margin{X} \margin{n+1}
\[ G_1 \hookrightarrow G_2 \hookrightarrow \ldots\ G_n \hookrightarrow G_{n\over{\unsure{-}}{+}1} \hookrightarrow \ldots \]
tel que $G_n$ est annulé par $p^n$ et $G_n(n-1) = G_{n-1}$ (on appelle un tel
système de groupes un \uline{système co-p-adique}) et tel que de plus $G_n$ soit
un $\Lambda_n$-module plat, alors $\lim\limits_{\rightarrow} G_n$ est un groupe p-divisible et de p-torsion.

4.2. \uline{Définition}. Un \uline{p-groupe de Barsotti-Tate sur $S$} (ou groupe
de Barsotti-Tate) est un groupe $G$ ($\mathcal{T}$-faisceau en groupes commutatifs)
p-divisible et de p-torsion tel que $G(1)$ soit un schéma fini localement
libre sur $S$.

Il est équivalent de dire que l'un quelconque des $G(n)$ est
un schéma fini localement libre sur $S$ (2.4).

\mspage{\hand{62}}

\clearpage
% Notes of the Montreal seminar, 1970 (cote 129, typed part): archive cote 129, site page 63 (PDF page 64).
% typescript. Draft compared: none.
% AI-transcribed by Claude Opus 5.5 (claude-opus-5-5, adaptive thinking, effort high), Message Batches API, batch mtl129_full (msgbatch_01AGnDKHwy6ATCM7ehUa6263), 2026-11-13; not yet checked by eye.
\scanpage{129/63}
\mspage{\hand{III-8}}

Si $G$ est un groupe de BT, alors, pour tout $n \geq 1$, $G(n)$
est un groupe de BT tronqué d'échelon $n$. Il se trouve que sur un
corps algébriquement clos, la réciproque est vraie: si $G'$ est un
groupe de BT tronqué d'échelon $n$, il existe un groupe de BT $G$ tel
que $G' = G(n)$ (la condition (*) de 3.2 est nécessaire si $n = 1$).

\margin{X} 5. \uline{\deleted{Suites}\inserted{Sorites} sur les groupes de Barsotti-Tate}

5.1. La catégorie des groupes de Barsotti-Tate peut se décrire sans
référence à la topologie \hand{$\mathcal{T}$} car elle est équivalente à la catégorie
des systèmes co-p-adiques $\{G_n\}$ où $G_n$ est un S-schéma en groupes fini
localement libre qui est un $\Lambda_n$-module plat.

La catégorie des groupes de BT est une catégorie additive.
Si $S'$ est un schéma sur $S$ et si $G$ est un groupe de BT sur $S$,
alors son image réciproque $G'$ sur $S'$ (restriction de $G$ aux
$S'$-schémas) est un groupe de BT sur $S'$. (On a en effet
$G' = \lim\limits_{\rightarrow} G(n) \times_S S'$ et les $G(n) \times_S S'$ sont des groupes de BT tronqués).

5.2. \uline{Proposition}. Soit $0 \rightarrow G' \rightarrow G \rightarrow G'' \rightarrow 0$ une suite exacte
de groupes; alors

a) Si $G'$ et $G''$ sont des groupes de BT, il en est de même de $G$;

b) Si $G'$ et $G$ sont des groupes de BT, il en est de même de $G''$;

a) Une extension de groupes p-divisibles (resp. de p-torsion) est
p-divisible (resp. de p-torsion) et la suite exacte induit (lemme du
Serpent) une suite exacte $0 \rightarrow G'(1) \rightarrow G(1) \rightarrow G''(1) \rightarrow G'/pG' \rightarrow \ldots$ mais

\mspage{\hand{63}}

\clearpage
% Notes of the Montreal seminar, 1970 (cote 129, typed part): archive cote 129, site page 64 (PDF page 65).
% typescript. Draft compared: none.
% AI-transcribed by Claude Opus 5.5 (claude-opus-5-5, adaptive thinking, effort high), Message Batches API, batch mtl129_full (msgbatch_01AGnDKHwy6ATCM7ehUa6263), 2026-11-13; not yet checked by eye.
\scanpage{129/64}
\mspage{\hand{III-9}}

comme $G'$ est p-divisible, $G(1)$ est une extension de groupes finis
localement libres et est donc de même nature.

b) Un quotient d'un groupe p-divisible (resp. de p-torsion) est
p-divisible (resp. de p-torsion). Comme précédemment, $G''(1)$ est
quotient de deux schémas en groupes finis localement libre, ce qui im-%
plique que $G''(1)$ l'est aussi (II, 1.6).

\uline{Remarque}. Si $G$ et $G''$ sont des groupes de BT, il est en général
faux que le noyau soit un groupe de BT. Par exemple, on a la suite
exacte $0 \to G(1) \to G \xrightarrow{p} G \to 0$ où $G(1)$ n'est pas p-divisible!

5.3. \uline{Rang (ou hauteur) d'un groupe de Barsotti-Tate}. $G(1)$ est un
groupe fini localement libre annulé par $p$. Pour tout $s \in S$, on a
donc
\[ \mathrm{rang}\ (G(1)_{s}) = p^{r(s)} \]
où la fonction $s \mapsto r(s)$ est localement constante. Des suites exactes
$0 \to G(1) \to G(n) \xrightarrow{p} G(n-1) \to 0$ on tire immédiatement que
\[ \mathrm{rang}\ (G(n)_{s}) = [\mathrm{rang}\ (G(1)_{s})]^{n} = p^{nr(s)} \]

L'entier $r(s)$ décrit donc le rang de tous les $G(n)$. Par
définition $r(s)$ s'appelle \uline{le rang} \uline{(ou la hauteur)} \uline{du groupe de Barsotti-}
\uline{Tate} $G$ \uline{au point} $s \in S$.

5.4. On a décrit la catégorie des groupes de BT sur $S$ à l'aide
de systèmes co-p-adiques. Signalons qu'on peut le faire à l'aide de
\uline{systèmes p-adiques:} Si $G$ est un groupe de Barsotti-Tate, on lui fait

\margin{64}

\clearpage
% Notes of the Montreal seminar, 1970 (cote 129, typed part): archive cote 129, site page 65 (PDF page 66).
% typescript. Draft compared: none.
% AI-transcribed by Claude Opus 5.5 (claude-opus-5-5, adaptive thinking, effort high), Message Batches API, batch pilotnext (msgbatch_01A9sdu8SYq23bEtqX9tUncQ), 2026-11-13; not yet checked by eye.
\scanpage{129/65}
\mspage{\hand{III-10}}

\margin{X} correspondre \deleted{G univoquement}\inserted{biunivoquement} le système projectif p-adique
\[ T_p(G) = (G(n)_{n\geq 1} \; ; \; G(n+1) \xrightarrow{p} G(n)) \]
où la multiplication par $p$ identifie $G(n)$ à $G(n+1)/p^nG(n+1)$ et
où $G(1)$ est annulé par $p$ et est représentable par un schéma fini
localement libre sur $S$.

5.5. \uline{Dual d'un groupe de Barsotti-Tate}. A tout groupe de BT
$G$ sur $S$, on associe donc le système p-adique $T_p(G)$. Ce système
projectif définit par dualité de Cartier un système co-p-adique de
groupes finis localement libres $G(n)^*$, annulés par $p^n$ et plats sur
$\Lambda_n$ (2.2). Par définition, le dual de $G$ est le groupe de BT $G^*$ as-%
socié à ce système co-p-adique:
\[ G^* = \lim_{\substack{\rightarrow \\ n}} G(n)^* \; . \]

On a en particulier $G^*(n) = G(n)^*$, et on obtient ainsi un
foncteur contravariant, compatible aux changements de base, qui définit
une équivalence entre la catégorie $BT(S)$ des p-groupes de Barsotti-%
Tate sur $S$ et son opposée, car on a un isomorphisme fonctoriel canonique
\[ G \xrightarrow{\sim} (G^*)^* \; . \]

5.6. \uline{Module de Dieudonné associé à un groupe de BT sur un corps}
\uline{parfait}. Reprenant les notations de II, 4.1 et désignant par $BT(k)$
la catégorie des p-groupes de BT sur le corps parfait $k$, on a le

\uline{Théorème} (cf. II, 5.3.1) Il existe une équivalence de catégories

\mspage{\hand{65}}

\clearpage
% Notes of the Montreal seminar, 1970 (cote 129, typed part): archive cote 129, site page 66 (PDF page 67).
% typescript. Draft compared: none.
% AI-transcribed by Claude Opus 5.5 (claude-opus-5-5, adaptive thinking, effort high), Message Batches API, batch mtl129_full (msgbatch_01AGnDKHwy6ATCM7ehUa6263), 2026-11-13; not yet checked by eye.
\scanpage{129/66}
\mspage{\hand{III-11}}

\[ D^{*}: [BT(k)]^{o} \xrightarrow{\ \approx\ } ( \hand{$\mathcal{D}$}\text{-modules, libres et de type fini sur } W) \]

Telle que
\[ [D^{*}(G)]^{\sigma} \simeq D^{*}(G^{(p)}) ; \quad D^{*}(\hand{F}_{G/k}) = F_{D^{*}(G)} ; \quad D^{*}(\hand{V}_{G}) = V_{D^{*}(G)} \]

Soit $G$ un groupe de BT sur un corps parfait. On pose
\[ M(n) = D^{*}(G(n)), \]
où $D^{*}$ est le foncteur ``classique'' défini en II, 4. Comme $G(n)$ est
annulé par $p^{n}$, $M(n)$ est annulé par $p^{n}$ et c'est donc un module sur
$W/p^{n}W = W_{n}$ .

Pour tout $i$, $1 \leq i \leq n-1$, les épimorphismes $G(n) \xrightarrow{p^{i}} G(n-i)$
définissent des monomorphismes
\[ M(n)/p^{n-i}M(n) = M(n-i) \xrightarrow{p^{i}} M(n) . \]
Ceci équivaut au fait que $M(n)$ est \uline{plat sur} $W_{n}$ (démonstration analogue
à celle de 2.2).

Enfin, le système co-p-adique des $G(n)$ définit un système
p-adique des modules $M(n)$: Les injections $G(n) \rightarrow G(n+1)$ définissent
des surjections
\[ M(n+1) \rightarrow M(n) \]
qui identifient $M(n)$ à $M(n+1)/p^{n+1}M(n)$.

Par définition, on pose:
\[ D^{*}(G) = \varprojlim M(n) \]

\mspage{\hand{66}}

\clearpage
% Notes of the Montreal seminar, 1970 (cote 129, typed part): archive cote 129, site page 67 (PDF page 68).
% typescript. Draft compared: none.
% AI-transcribed by Claude Opus 5.5 (claude-opus-5-5, adaptive thinking, effort high), Message Batches API, batch mtl129_full (msgbatch_01AGnDKHwy6ATCM7ehUa6263), 2026-11-13; not yet checked by eye.
\scanpage{129/67}
\mspage{\hand{III-12}}

$D^*(G)$ est un module sur $W = \underset{\leftarrow}{\lim}\, W_n$; comme limite p-adique
de $W_n$-modules plats, il est plat sur l'anneau local $W$ et est donc
\uline{libre}. Les $M(n)$ étant de longueur finie, $D^*(G)$ est de \uline{type fini}.
De plus, $D^*(G)$ est muni de morphismes semi-linéaires $F$ et $V$ induits
par les morphismes correspondants sur $M(n)$.

6. \uline{Exemples et types particuliers de groupes de Barsotti-Tate}

6.1. Soit $A$ un schéma abélien sur $S$ (i.e. un schéma en groupes
commutatifs, propre et lisse sur $S$, et à fibres connexes). Un schéma
abélien est p-divisible: c'est bien connu dans le cas où la base est
le spectre d'un corps algébriquement clos (cf. S. Lang, ``Abelian
Varieties'') et, si $S$ est quelconque, la surjectivité de $\mathrm{pId}_A$ résulte
de ce cas particulier grâce au critère de platitude fibre par fibre
(EGA, IV, 11.3.10).

Si $d$ désigne la dimension relative de $A$ sur $S$, on peut
montrer que
\[ \text{rang } A(1) = p^{2d} . \]

Le groupe $A(1)$ est donc fini; il est évidemment de présen-%
tation finie et, étant plat, il est localement libre. On en déduit que
$A(\infty) = \underset{\overrightarrow{n}}{\lim}\, A(n)$ est un groupe de BT de rang $2d$. C'est le \uline{p-groupe
de BT associé au schéma abélien $A$}, et cette construction peut se faire
pour tout nombre premier $p$.

Si toutes les caractéristiques résiduelles de $S$ sont diffé-%

\margin{\hand{67}}

\clearpage
% Notes of the Montreal seminar, 1970 (cote 129, typed part): archive cote 129, site page 68 (PDF page 69).
% typescript. Draft compared: none.
% AI-transcribed by Claude Opus 5.5 (claude-opus-5-5, adaptive thinking, effort high), Message Batches API, batch mtl129_full (msgbatch_01AGnDKHwy6ATCM7ehUa6263), 2026-11-13; not yet checked by eye.
\scanpage{129/68}
\mspage{\hand{III - 13}}

rentes de $p$, on montre que $A(1)$ est étale sur $S$, et on dira alors
que $A(\infty)$ est ind-étale.

Notons par ailleurs qu'il existe une notion de dualité pour
les variétés abéliennes. Si $A^*$ désigne la variété duale de la variété
abélienne $A$, on montre aisément que son groupe de BT associé $A^*(\infty)$
s'identifie au dual, en tant que groupe de BT de $A(\infty)$.

6.2. Par définition, on dit qu'un groupe de BT sur $S$ est \uline{un groupe}
\uline{de Barsotti-Tate ind-étale} si $G(1)$ est étale. Il est équivalent de dire
(les $G(n)$ admettant une suite de composition en groupes isomorphes à
$G(1)$) que tous les $G(n)$ sont étales.

Le foncteur $G \mapsto T_p(G)$ (5.4) définit une équivalence de caté-%
gories entre la catégorie des groupes de BT ind-étales et la catégorie
des \uline{faisceaux p-adiques constants tordus sans torsion} (i.e. systèmes
p-adiques de groupes finis localement libres $G_n$, étales sur $S$ et plats
sur $\Lambda_n$).

Si $S$ est connexe, la catégorie des faisceaux p-adiques cons-%
tants tordus sans torsion est bien connue: si l'on choisit un point géo-%
métrique $\overline{s} \in S$, on définit une équivalence de cette catégorie avec la
catégorie des $Z_p$-modules libres de type fini sur lesquels le groupe fon-%
damental $\pi_1(S,\overline{s})$ opère continuement en associant au système $(G_n)_{n\geq 1}$
le $Z_p$-module $\varprojlim\, (G_n)_{\overline{s}}$ , $\pi_1(S,\overline{s})$ opérant par monodromie sur $(G_n)_{\overline{s}}$ .

6.3. \uline{Exemple.} On dit qu'un schéma en groupes $T$ sur $S$ est un
\uline{tore} s'il est, localement pour la topologie étale, isomorphe à $\mathbb{G}^r_{m,S}$ .

\margin{68}

\clearpage
% Notes of the Montreal seminar, 1970 (cote 129, typed part): archive cote 129, site page 69 (PDF page 70).
% typescript. Draft compared: none.
% AI-transcribed by Claude Opus 5.5 (claude-opus-5-5, adaptive thinking, effort high), Message Batches API, batch mtl129_full (msgbatch_01AGnDKHwy6ATCM7ehUa6263), 2026-11-13; not yet checked by eye.
\scanpage{129/69}
\mspage{\hand{III-14}}

Un tore  T  est p-divisible et  T(1)  est un groupe fini localement libre
de rang  $p^{r}$,  r  étant la dimension relative de  T  sur  S  (ces proprié-%
tés sont locales et sont vérifiées par  $\mathbb{G}^{r}_{m,S}$).  On en déduit que  $T(\infty)$
est un groupe de  BT  de rang  r.

Comme exemple particulier important, on a le groupe de  BT
\[ \hand{\mu} = \mathbb{G}_{m,S}(\infty) = \varinjlim_{n} \hand{\mu}_{p^{n},S} \ . \]
\margin{$\times$ \ $\mu$}

6.4.  \uline{Exemple.}  Soit  G  un schéma en groupes sur  S,  extension d'un
schéma abélien  A  par un tore  T.  Alors  G  est p-divisible,  G(1)  est
fini localement libre sur  S  et  $G(\infty)$  est un groupe de  BT  de rang
2d\hand{+}r  (d  et  r  désignant respectivement les dimensions relatives de  A
et  T  sur  S).  De plus, on a une suite exacte
\[ 0 \to T(\infty) \to G(\infty) \to A(\infty) \to 0 \ . \]

6.5.  Comme y incite l'exemple (6.3), on dira qu'un groupe de  BT
G  est \uline{toroïdal} si  G(1)  est de type multiplicatif.  Il est équivalent
de dire que tous les  G(n)  sont de type multiplicatif, ou encore que le
dual  $G^{*}$  (5.5) est ind-étale.

Le foncteur  $G \mapsto T_{p}(G^{*})$  définit une anti-équivalence entre la
catégorie des groupes de  BT  toroidaux et celle des faisceaux p-adiques
constants tordus sans torsion.  Composant ce foncteur avec la dualité
dans les faisceaux p-adiques constants tordus sans torsion
\[ T_{p}(G^{*}) \mapsto \underline{\mathrm{Hom}}_{gr}(G^{*},\mathbb{Z}_{p}) \overset{\text{def}}{=} (\underline{\mathrm{Hom}}_{gr}(G(n)^{*},\Lambda_{n}))_{n\geq 1} \ , \]

\mspage{\hand{69}}

\clearpage
% Notes of the Montreal seminar, 1970 (cote 129, typed part): archive cote 129, site page 70 (PDF page 71).
% typescript. Draft compared: none.
% AI-transcribed by Claude Opus 5.5 (claude-opus-5-5, adaptive thinking, effort high), Message Batches API, batch mtl129_full (msgbatch_01AGnDKHwy6ATCM7ehUa6263), 2026-11-13; not yet checked by eye.
\scanpage{129/70}
\mspage{\hand{III-15}}

on obtient une équivalence au lieu d'une anti-équivalence. Cette équi-%
valence peut s'expliciter comme suit:
\[
G \mapsto \underline{\mathrm{Hom}}_{gr}(\hand{\mu},G) \overset{\mathrm{def}}{=\!=} (\underline{\mathrm{Hom}}_{gr}(\hand{\mu}(n),G(n)))_{n\geq 1}
\]

6.6. On dit qu'un groupe de $BT$ $G$ sur $S$ est \uline{à fibres connexes}
si $G(1)$ est infinitésimal. Il revient au même de dire que tous les
$G(n)$ sont infinitésimaux.

Si toutes les caractéristiques résiduelles de $S$ sont égales
à $p$, les groupes de $BT$ toroidaux sont à fibres connexes (II,2.2).

On montrera plus loin que, si $p$ est localement nilpotent sur
$S$, un groupe de $BT$ est à fibres connexes si et seulement si il est
aussi un groupe de Lie formel.

Nous allons donner ici une première définition élémentaire des
groupes de Lie formels et étudier à quelles conditions un groupe de Lie
formel est un groupe de $BT$.

6.7. On appelle \uline{variété de Lie formelle ponctuée} (ou \uline{augmentée}) sur
$S$ un faisceau $X$ sur $S$ (pour la topologie $\hand{\mathcal{T}}$, (f.p.p.f.) $\leq \hand{\mathcal{T}} \leq$ (f.p.q.c.))
muni d'une section $S \overset{\hand{\xi}}{\to} X$ et tel que localement (pour la topologie de
Zariski) sur $X$, on ait
\[
X \simeq \varinjlim_{n} \operatorname{Spec} \underline{\hand{\mathcal{O}}}_S[[T_1,\ldots,T_r]]/(T_1,\ldots,T_r)^n .
\]

\uline{Remarques}. 1) On peut expliciter ce dernier isomorphisme en disant
que pour tout $S' \to S$, $S'$ quasi-compact, on a

\mspage{\hand{70}}

\clearpage
% Notes of the Montreal seminar, 1970 (cote 129, typed part): archive cote 129, site page 71 (PDF page 72).
% typescript. Draft compared: none.
% AI-transcribed by Claude Opus 5.5 (claude-opus-5-5, adaptive thinking, effort high), Message Batches API, batch mtl129_full (msgbatch_01AGnDKHwy6ATCM7ehUa6263), 2026-11-13; not yet checked by eye.
\scanpage{129/71}
\mspage{\hand{III-1\unsure{6}}}

\[
X(S') = \varinjlim_{n} \mathrm{Hom}_S\ (S',\ \mathrm{Spec}\ \underline{\mathcal{O}}_S[[T_1,\ldots,T_r]]/(T_1,\ldots,T_r)^n)
\]
\[
= \varinjlim_{n} \{f_1^*,f_2,\ldots,f_r \in \Gamma(S',\underline{\mathcal{O}}_S)\ \mid\ f_i^n = 0,\ \ 1 \leq i \leq r\}
\]

2) $X$ est en particulier limite inductive de schémas finis
et radiciels au dessus de $S$.

On appelle \uline{groupe de Lie formel sur $S$} un faisceau en groupes
sur $S$ tel que sa section unité lui donne une structure de variété de
Lie formelle. Comme le produit de deux variétés de Lie formelles est
encore une variété de Lie formelle, il est équivalent de dire qu'un grou-%
pe de Lie formel est un objet en groupes de la catégorie des variétés de
Lie formelles sur $S$.

Dans ce qui suit, tous les groupes de Lie formels seront suppo-%
sés commutatifs.

Soit $G$ un groupe de Lie formel sur $S$. On vérifie aisément
que, si $p$ est localement nilpotent sur $S$, $G$ est de $p$-torsion, c'est-%
à-dire que
\[
G = G(\infty) = \varinjlim_{n} G(n) \ ,
\]
de sorte que $G$ est un groupe de Barsotti-Tate si et seulement si les
deux conditions suivantes sont vérifiées:

1) $G$ est $p$-divisible.

2) $G(1)$ est représentable par un schéma fini localement libre.

Il se trouve que \uline{si $S$ est artinien local}, d'unique point $s \in S$,

\mspage{\hand{71}}

\clearpage
% Notes of the Montreal seminar, 1970 (cote 129, typed part): archive cote 129, site page 72 (PDF page 73).
% typescript. Draft compared: none.
% AI-transcribed by Claude Opus 5.5 (claude-opus-5-5, adaptive thinking, effort high), Message Batches API, batch mtl129_full (msgbatch_01AGnDKHwy6ATCM7ehUa6263), 2026-11-13; not yet checked by eye.
\scanpage{129/72}
\mspage{\uline{III}-17}

ces deux conditions sont équivalentes et équivalent à chacune des con-%
ditions (ne faisant intervenir que la fibre géométrique de $G$):

3) $G_{\bar{s}}$ est p-divisible.

4) Il est impossible de plonger le groupe formel additif $\overline{\mathbb{G}}_a$
(complété formel de $\mathbb{G}_a$ le long de sa section unité) dans $G_{\bar{s}}$ .

Dans le cas général, on pense que les conditions 1) et 2) sont
encore équivalentes; cela résulterait de la conjecture suivante, vraie
dans le cas artinien local (SGA 3, VII, \quad ):

Si $u\colon G \to G'$ est un S-morphisme de variétés de Lie formelles
de même dimension relative sur $S$, on a les équivalences:

- $u$ est un épimorphisme

- Ker $(u)$ est représentable par un schéma fini

- Ker $(u)$ est représentable par un schéma fini localement libre.

Par contre, le fait que le groupe de Lie formel $G$ soit un
groupe de BT ne se vérifie pas en général sur les fibres, comme le mon-%
tre l'exemple du groupe de Lie formel associé à la courbe elliptique mo-%
dulaire $A$ sur une courbe $S$ de caractéristique $p$, puisqu'il y a des
points $s \in S$ tels que l'invariant de Hasse de $A_s$ est nul.

Cet exemple suggère cependant la conjecture suivante: Pour
qu'un groupe de Lie formel $G$ sur $S$ soit de BT, il faut et il suffit
que ses fibres $G_s$ soient des groupes de BT et que le rang de $G_s$
(en tant que BT) soit une fonction localement constante en $s \in S$.

\mspage{\hand{72}}

\clearpage
% Notes of the Montreal seminar, 1970 (cote 129, typed part): archive cote 129, site page 73 (PDF page 74).
% typescript. Draft compared: none.
% AI-transcribed by Claude Opus 5.5 (claude-opus-5-5, adaptive thinking, effort high), Message Batches API, batch mtl129_full (msgbatch_01AGnDKHwy6ATCM7ehUa6263), 2026-11-13; not yet checked by eye.
\scanpage{129/73}
\mspage{\hand{III - 18}}

7. \uline{Suite de composition d'un groupe de Barsotti-Tate}

7.1. Supposons d'abord $S$ réduit à un point $s$. Si $G$ est un
groupe de Barsotti-Tate sur $S$, on pose par définition
\[ G^{o} = \underset{\vec{n}}{\lim}\ G(n)^{o} \qquad\qquad (G(n)^{o} = \text{composante neutre de } G(n)) \]
\[ G^{et} = \underset{\vec{n}}{\lim}\ G(n)/G(n)^{o} \ , \]
et l'on obtient une suite exacte
\[ 0 \to G^{o} \to G \to G^{et} \to 0 \]
On vérifie immédiatement que $G^{o}$ est un groupe de BT, d'où l'on déduit
(5.2) que $G^{et}$ l'est aussi, de sorte que l'on obtient $G$ comme extension
d'un groupe de BT ind-étale par un groupe de BT à fibre connexe. Cet-%
te décomposition est évidemment canonique et fonctorielle en $G$.

7.2. Si la caractéristique résiduelle de $S$ est $p$ et si l'on pose
\[ G^{tor} = \underset{\vec{n}}{\lim}\ G(n)^{mult} \]
où $G(n)^{mult}$ désigne le plus grand sous-groupe de type multiplicatif de
$G(n)$, $G(n)^{mult} \subset G(n)^{o}$, on obtient une filtration
\[ \{0\} \subset G^{tor} \subset G^{o} \subset G \]

Il est immédiat que $G^{tor}$ est un groupe de BT; c'est le plus
grand sous-groupe de BT toroidal de $G$.

On en déduit que $G^{o}/G^{tor}$ est un groupe de BT; c'est un

\margin{73}

\clearpage
% Notes of the Montreal seminar, 1970 (cote 129, typed part): archive cote 129, site page 74 (PDF page 75).
% typescript. Draft compared: none.
% AI-transcribed by Claude Opus 5.5 (claude-opus-5-5, adaptive thinking, effort high), Message Batches API, batch mtl129_full (msgbatch_01AGnDKHwy6ATCM7ehUa6263), 2026-11-13; not yet checked by eye.
\scanpage{129/74}
\mspage{\hand{\unsure{III-17|III-19}}}

groupe à fibre connexe et ind-unipotent (i.e. $G^{o}/G^{tor}(1) = G(1)^{o}/G(1)^{mult}$
est unipotent). Cette suite de composition canonique est fonctorielle
en $G$, commute aux changements de base et par passage au dual, on véri-%
fie que, par construction même, on a des isomorphismes canoniques
\[ (G/G^{o})^{*} \simeq (G^{*})^{tor} , \quad (G/G^{tor})^{*} \simeq (G^{*})^{o} . \]

7.3. Si la base $S$ est quelconque, on en peut pas en général écrire
un groupe de BT $G$ comme extension d'un groupe de BT ind-étale $G''$
par un groupe de BT à fibres connexes $G'$.

En effet, si on a une telle extension, on sait que l'on obtient
alors une suite exacte
\[ 0 \to G'(1) \to G(1) \to G''(1) \to 0 . \]

Pour tout $s \in S$, considérons alors le rang séparable de
$G(1)_{s}$. On a:

rang sep $(G(1)_{s})$ = card $G(1)(\overline{s})$ = rang sep $(G'(1)_{s}) \times$ rang sep $(G''(1)_{s})$.

Mais le rang séparable de $G'(1)_{s}$ est $1$ car $G'(1)$ est infi-%
nitésimal et le rang séparable de $G''(1)_{s}$ est une fonction localement cons-%
tante de $s \in S$ car $G''(1)$ est étale sur $S$, de sorte que l'on en dé-%
duit que le rang séparable de $G(1)_{s}$ est une fonction localement constante
de $s$.

Or cela n'est pas vrai en général, comme le montre l'exemple de
la variété elliptique modulaire $A$ sur une courbe $S$ en caractéristique

\mspage{\hand{74}}

\clearpage
% Notes of the Montreal seminar, 1970 (cote 129, typed part): archive cote 129, site page 75 (PDF page 76).
% typescript. Draft compared: none.
% AI-transcribed by Claude Opus 5.5 (claude-opus-5-5, adaptive thinking, effort high), Message Batches API, batch mtl129_full (msgbatch_01AGnDKHwy6ATCM7ehUa6263), 2026-11-13; not yet checked by eye.
\scanpage{129/75}
\mspage{\hand{III-20}}

p, où l'on a des points isolés $s$ où l'invariant de Hasse de $A_s$ est
nul. On a cependant le résultat suivant:

7.4. \uline{Proposition}. Soit $G$ un groupe de Barsotti-Tate sur $S$. Les
propriétés suivantes sont équivalentes:

i) $G$ est extension d'un groupe de BT ind-étale par un groupe
de BT à fibres connexes.

ii) Pour tout $n \in \mathbb{N}$, $G(n)$ est extension d'un groupe fini
étale par un groupe fini infinitésimal.

ii)$^{\text{bis}}$ $G(1)$ est extension d'un groupe fini étale par un groupe
fini infinitésimal.

iii) Pour tout $n \in \mathbb{N}$, la fonction $s \mapsto \text{rang sep } G(n)_s$ est lo-%
calement constante.

iii)$^{\text{bis}}$ La fonction $n \mapsto \text{rang sep } G(1)_s$ est localement constante.

L'équivalence de (iii) et (iii)$^{\text{bis}}$ résulte de ce que le rang sé-%
parable de $G(n)_s$ est la puissance n-ième du rang séparable de $G(\unsure{i|1})_s$ .
D'autre part, (i) résulte immédiatement de (ii), et, d'après ce qui précède,
la seule implication à démontrer est (iii) $\Rightarrow$ (ii). On utilise pour cela
le lemme suivant (EGA, IV, ):

\uline{Lemme}. Soit $f : X \to S$ un morphisme fini localement libre de schémas.
Pour que $f$ se factorise en $f'$ étale et $f''$ radiciel et surjectif, il
\begin{diagramblock}\begin{tikzcd}[column sep=small]
X \arrow[rr, "f''"] \arrow[dr, "f"'] & & Y \arrow[dl, "f'"] \\
& S &
\end{tikzcd}\end{diagramblock}
faut et il suffit que la fonction $s \mapsto \text{rang sep } (X_s)$
soit localement constante. La factorisation est alors
unique et fonctorielle en $f$.

(De la fonctorialité, il résulte que si $X$ est un schéma en groupes sur
$S$, $Y$ l'est aussi.)

\margin{75}

\end{document}
